\documentclass[11pt,a4paper]{article}
\usepackage[top=2.5cm,bottom=2.8cm,left=2.5cm,right=2.5cm]{geometry}
\usepackage{amsmath,amssymb,amsthm}
\usepackage{hyperref,xcolor,enumitem,booktabs}

\hypersetup{colorlinks=true,linkcolor={blue!60!black},
  citecolor={blue!60!black},urlcolor={blue!60!black}}

\newtheorem{theorem}{Theorem}[section]
\newtheorem{lemma}[theorem]{Lemma}
\newtheorem{proposition}[theorem]{Proposition}
\newtheorem{corollary}[theorem]{Corollary}
\theoremstyle{definition}\newtheorem{definition}[theorem]{Definition}
\theoremstyle{remark}
\newtheorem{remark}[theorem]{Remark}
\newtheorem{conjecture}[theorem]{Conjecture}

\newcommand{\vp}{\varphi}
\newcommand{\lam}{\lambda}
\newcommand{\Ja}{J_{2a}}
\newcommand{\Jfour}{J_4}
\newcommand{\ms}{\mu^*}
\newcommand{\TCMB}{T_{\mathrm{CMB}}}
\newcommand{\Lcond}{\Lambda_{\mathrm{cond}}}
\newcommand{\rCMB}{\rho_{\mathrm{CMB}}}
\newcommand{\SU}{\mathrm{SU}}
\newcommand{\Efb}{E_{\mathrm{fb}}}
\newcommand{\Kfb}{K_{\mathrm{fb}}}
\newcommand{\twoI}{2I}
\newcommand{\twoT}{2T}
\newcommand{\QH}{Q_H}

\title{\textbf{The $\QH=3$ Sector of the Density-Feedback Hopfion}}
\author{Frederick Manfredi}
\date{2026}

\begin{document}
\maketitle

\begin{abstract}
We extend the density-feedback Faddeev--Niemi framework of
Papers~I--XIV to the $\QH=3$ topological sector,
whose saddle-point soliton has a trefoil-knotted flux tube.
Six results are established analytically.
(i) Via the McKay correspondence
$\twoT\leftrightarrow E_6\leftrightarrow\SU(3)_1$,
the WZW model for the $\QH=3$ sector is $\SU(3)_1$.
(ii)~The T-matrix closure condition
$\cos(2\pi Q T_j)=0$ for all $\SU(3)_1$ primaries
is satisfied at the unique minimum
$Q_{\mathrm{group}}^{(3)}=3$.
(iii)~The conformal weight $h_{(1,0)}=\tfrac{1}{3}$
equals the fractional electric charge of the down quark exactly.
(iv)~The virial fixed point $V^{*\,(3)}=\vp$ is established by
the IVT argument of Paper~III, giving $\Kfb^{(3)}=2\vp\Ja^{(3)}$
and $\Efb^{(3)}=2N_3^2/\vp^{17}$.
(v)~Using the Frenet--Serret transformation, we prove that torsion
vanishes \emph{at all orders} by azimuthal averaging for the
symmetric profile $f_0$ (Theorem~\ref{P15:thm:torsion_cancels}),
leaving a leading-order formula
$J_4^{(3)}/J_{2a}^{(3)}=3/(4C_3^{*2})+\langle\kappa^2\rangle/2$,
with an $O(\kappa^2\tau^2)$ next-order correction identified in
Remark~\ref{P15:rem:torsion_scope}.
We further identify the three crossing regions of the trefoil
as the $\QH=3$ analogue of the torus inner wall: the loci of
enhanced condensate density, with $\mathbb{Z}_3$ symmetry,
that constitute the spatial imprint of the three quark colours.
Colour confinement follows topologically from the profile
boundary conditions already present in Paper~I~\cite{Paper1}.
(vi)~The Bogomolny parameter $\lam_3=\vp^6$ is proved
(Theorem~\ref{P15:thm:sector_assignment} and Corollary~\ref{P15:cor:lambda3_proved}):
the triple-fusion decomposition
$j_{1/2}^{\otimes 3}=j_{1/2}^{(\times 2)}\oplus j_{3/2}^{(\times 1)}$
at $k=3$, combined with the topological non-degeneracy
$\pi_3(S^2)=\mathbb{Z}$, uniquely forces $\QH=3$ into
the $j=3/2$ simple-current sector, supported by seven of eight
independent criteria in Table~\ref{P15:tab:sector_evidence}.
The normalisation identity $\vp^9\sqrt{\Ja^{(3)}\Jfour^{(3)}}=3$
(Theorem~\ref{P15:thm:norm3}) and condensate volume $V_3=9/(5\vp^{17})$
(Corollary~\ref{P15:cor:V3}) follow unconditionally as derived consequences.
The $O(\langle\kappa^4\rangle)$ curvature correction is computed in
closed form via Beta-function integration, giving $C_3^*=2.5062$.
The remaining problem is numerical: a topology-preserving
saddle finder to confirm $C_3^*=2.5062$ directly from the $\QH=3$
Euler--Lagrange profile, which does not affect the main results.
\end{abstract}

\noindent\textbf{Keywords:}
Faddeev--Niemi Hopfion; trefoil knot; $\QH=3$; $\SU(3)_1$ WZW;
quark charge; T-matrix closure; Frenet--Serret; colour confinement;
Shafranov shift; McKay correspondence.

\tableofcontents\bigskip

%══════════════════════════════════════════════════════════════════════════════
\section{Introduction}
\label{P15:sec:intro}
%══════════════════════════════════════════════════════════════════════════════

Papers~I--XIV established that the density-feedback
Faddeev--Niemi Hopfion in the $\QH=2$ sector reproduces, from
the single input $\TCMB$, the electron mass, fine-structure
constant, electroweak scale, and Higgs boson mass.
The central identity of Paper~XII is
\begin{equation}
  \vp^9\sqrt{\Ja\Jfour} = 10 = \mathrm{ord}(q_{\twoI}),
  \label{P15:eq:norm2}
\end{equation}
and the Bogomolny parameter $\lam=\vp^6$ is uniquely fixed by
the matching condition $Q_{\mathrm{num}}(\lam)=10$ proved in
Paper~XIV~\cite{Paper14}.

The $\QH=3$ topological sector is the subject of this paper.
In the pure Faddeev--Niemi model~\cite{Faddeev1997}
the $\QH=3$ soliton is a trefoil-knotted flux tube with
$\mathbb{Z}_3$ symmetry~\cite{Battye1998}.
The binary tetrahedral group $\twoT\subset\SU(2)$ (order~24)
sits at the $E_6$ node of the ADE McKay chain:
\begin{equation}
  \twoI\;\leftrightarrow\;E_8\;\leftrightarrow\;\SU(2)_3
  \quad\text{(lepton sector)},
  \qquad
  \twoT\;\leftrightarrow\;E_6\;\leftrightarrow\;\SU(3)_1
  \quad\text{(quark sector)}.
  \label{P15:eq:mckay}
\end{equation}
The $\SU(3)_1$ WZW model carries all the structure required for colour:
three primary fields, fractional conformal weight $\tfrac{1}{3}$,
and $\mathbb{Z}_3$ outer automorphism matching the three quark colours.

\paragraph{Main results.}
\begin{enumerate}[noitemsep,label=(\arabic*)]
\item $Q_{\mathrm{group}}^{(3)}=3$
  (Theorem~\ref{P15:thm:Qgroup3}, exact arithmetic).
\item $h_{(1,0)}=\tfrac{1}{3}$: quark charge derived from the
  trinification electric-charge generator
  (Theorem~\ref{P15:thm:charge}, Theorem~\ref{P15:thm:charge_generator},
  Corollary~\ref{P15:cor:charge_eq_weight}).
\item $V^{*\,(3)}=\vp$ and $\Kfb^{(3)}=2\vp\Ja^{(3)}$
  (Theorem~\ref{P15:thm:virial}, exact algebra using Paper~III IVT).
\item Torsion vanishes \emph{at all orders} by azimuthal averaging
  for the symmetric profile $f_0$ (Theorem~\ref{P15:thm:torsion_cancels});
  a small $O(\kappa^2\tau^2)$ correction at next order is
  identified in Remark~\ref{P15:rem:torsion_scope};
  leading-order $J_4^{(3)}/J_{2a}^{(3)}$ formula
  (equation~\eqref{P15:eq:r3formula}).
\item Crossing regions as the trefoil inner wall; $\mathbb{Z}_3$
  colour imprint; confinement from profile topology
  (Section~\ref{P15:sec:geometry}).
\item $\twoT\subset\twoI$; three quark colours as $\omega$-twisted
  fermionic irreps of $\twoT$; leading-order $m_u=m_d$
  (Theorem~\ref{P15:thm:subgroup}, Proposition~\ref{P15:prop:colours},
  Theorem~\ref{P15:thm:massdegen}).
\item $\lam_3=\vp^6$ (Theorem~\ref{P15:thm:sector_assignment},
  Corollary~\ref{P15:cor:lambda3_proved}),
  proved via the triple-fusion decomposition
  $j_{1/2}^{\otimes 3}=j_{1/2}^{(\times 2)}\oplus j_{3/2}^{(\times 1)}$
  at $k=3$ plus topological non-degeneracy $\pi_3(S^2)=\mathbb{Z}$,
  supported by seven of eight independent criteria
  (Table~\ref{P15:tab:sector_evidence}).
  The normalisation identity $\vp^9\sqrt{\Ja^{(3)}\Jfour^{(3)}}=3$
  (Theorem~\ref{P15:thm:norm3}) and condensate volume $V_3=9/(5\vp^{17})$
  (Corollary~\ref{P15:cor:V3}) follow as derived consequences.
\item The $O(\langle\kappa^4\rangle)$ curvature correction is computed
  in closed form via Beta-function integration, giving
  $C_3^*=2.5062$ (Proposition~\ref{P15:prop:kappa4}); the
  $2I\to 2T$ branching rules are proved by character theory
  (Proposition~\ref{P15:prop:branching}), with the 24-cell geometry
  of $\twoT$ explaining the integer normalisation
  (Remark~\ref{P15:rem:F4}).
\end{enumerate}

%══════════════════════════════════════════════════════════════════════════════
\section{The $(E_6)_1$ and SU(3)$_1$ WZW Models for the Q$_H$=3 Sector}
\label{P15:sec:wzw}
%══════════════════════════════════════════════════════════════════════════════

\subsection{McKay correspondence and the correct WZW chain}

The McKay correspondence associates to each finite subgroup
$G\subset\SU(2)$ an affine Dynkin diagram whose nodes are the
irreducible representations of $G$.
For $\twoT$ (order~24, seven irreps of dimensions $1,1,1,2,2,2,3$),
the McKay graph is the \emph{affine} $\widehat{E}_6$ diagram
(seven nodes).
Via the McKay correspondence, the associated WZW model is
$(E_6)_1$ (affine $E_6$ at level~1), with central charge
\begin{equation}
  c_{(E_6)_1} = \frac{k\cdot\dim E_6}{k+h^\vee}\bigg|_{k=1,\,h^\vee=12}
  = \frac{78}{13} = 6.
  \label{P15:eq:cE6}
\end{equation}
The $(E_6)_1$ model has three primary fields:
the vacuum~$\mathbf{1}$, and the two 27-dimensional
representations $\mathbf{27}$ and $\overline{\mathbf{27}}$,
with conformal weights $h_{(0)}=0$ and
\begin{equation}
  h_{\mathbf{27}} = h_{\overline{\mathbf{27}}} = \tfrac{2}{3}.
  \label{P15:eq:h27}
\end{equation}

\begin{remark}[Clarification on the McKay chain]
The $\widehat{E}_6$ diagram (7~nodes) and the affine
$\widehat{A}_2=\widehat{\mathfrak{sl}(3)}$ diagram (3~nodes,
underlying $\SU(3)_1$) are \emph{different} diagrams.
The WZW model directly associated to $\twoT$ via the McKay
correspondence is $(E_6)_1$, not $\SU(3)_1$.
The connection to $\SU(3)_1$ is established through a
conformal embedding (Section~\ref{P15:sec:embed}), which is
the $E_6\supset\SU(3)^3$ branching at level~1.
The same structural convention appears in Papers~I--XIV,
where the McKay graph of $\twoI$ is $\widehat{E}_8$ (9~nodes,
giving $(E_8)_1$ with $c=8$), while $\SU(2)_3$ ($c=9/5$)
is identified via the Pentagon number $k+2=5$ of Paper~V
rather than through a conformal embedding of $(E_8)_1$.
The $\QH=3$ sector improves on this: the bridge to $\SU(3)_1$
is an explicit conformal embedding with exact central-charge match.
\end{remark}

\subsection{Conformal embedding $(E_6)_1\supset\SU(3)_1^{\times 3}$}
\label{P15:sec:embed}

The Lie algebra $E_6$ contains $\mathfrak{su}(3)^3$ as a
maximal subalgebra via the branching
$E_6\supset\SU(3)_L\times\SU(3)_R\times\SU(3)_C$,
with the $\mathbf{27}$ decomposing as
\begin{equation}
  \mathbf{27} \;\to\;
  (\mathbf{3},\bar{\mathbf{3}},\mathbf{1})
  \oplus(\bar{\mathbf{3}},\mathbf{1},\mathbf{3})
  \oplus(\mathbf{1},\mathbf{3},\bar{\mathbf{3}}).
  \label{P15:eq:27branch}
\end{equation}
At level~1 this gives a \emph{conformal embedding}:
\begin{equation}
  (E_6)_1 \supset \SU(3)_1 \times \SU(3)_1 \times \SU(3)_1,
  \label{P15:eq:confemb}
\end{equation}
verified by the central-charge identity
\begin{equation}
  c_{(E_6)_1} = 6 = 3\times c_{\SU(3)_1} = 3\times 2.
  \label{P15:eq:cematch}
\end{equation}
The conformal weight of the $(\mathbf{3},\bar{\mathbf{3}},\mathbf{1})$
component under $\SU(3)_1^{\times 3}$ is
$h_{\mathbf{3}}+h_{\bar{\mathbf{3}}}+0=\tfrac{1}{3}+\tfrac{1}{3}
=\tfrac{2}{3}$,
matching $h_{\mathbf{27}}=\tfrac{2}{3}$ exactly.

The three $\SU(3)_1$ factors in~\eqref{P15:eq:confemb} are isomorphic
copies with identical conformal data ($c=2$, $h=0,\tfrac13,\tfrac13$
for the three primaries); they are permuted by the $\mathbb{Z}_3$
outer automorphism of $E_6$ and identified with the three quark
colours in Section~\ref{P15:sec:2T2I}. For the electric-charge
derivation of Section~\ref{P15:sec:quarkcharge}, the relevant factor is
the electroweak $\SU(3)_L\subset(E_6)_1$ of the trinification
$E_6\supset\SU(3)_C\times\SU(3)_L\times\SU(3)_R$: the WZW data
$h_{(1,0)}=\tfrac13$ is identical across all three factors, but
only $\SU(3)_L$ (or, equivalently, $\SU(3)_R$) carries the Standard
Model electric-charge generator (Theorem~\ref{P15:thm:charge_generator}).
Colour-labelling (which leg is ``colour 1/2/3'') and
charge-generation (which leg's Cartan contains $Q$) are therefore
two independent roles of the same conformal embedding, not in
tension with each other.

\subsection{Primary fields and conformal data}

Two complementary descriptions coexist:

\paragraph{$(E_6)_1$ level.}
Three primaries $\mathbf{1},\mathbf{27},\overline{\mathbf{27}}$
with $c=6$ and
\begin{equation}
  T_{(0)} = -\tfrac{1}{4},\qquad
  T_{\mathbf{27}} = T_{\overline{\mathbf{27}}} = \tfrac{5}{12}.
  \label{P15:eq:TE6}
\end{equation}

\paragraph{$\SU(3)_1$ level ($\SU(3)_L$ electroweak factor, used for the charge computation of Section~\ref{P15:sec:quarkcharge}).}
Three primaries $(0,0),(1,0),(0,1)$ with $c=2$ and
\begin{equation}
  T_{(0,0)} = -\tfrac{1}{12},\qquad
  T_{(1,0)} = T_{(0,1)} = \tfrac{1}{4}.
  \label{P15:eq:TSU3}
\end{equation}
The two descriptions are related by~\eqref{P15:eq:confemb}:
the $\SU(3)_1$ T-matrix exponents are embedded within the
$(E_6)_1$ spectrum.

%══════════════════════════════════════════════════════════════════════════════
\section{T-Matrix Closure: $Q_{\mathrm{group}}^{(3)}=3$ at Both Levels}
\label{P15:sec:qgroup}
%══════════════════════════════════════════════════════════════════════════════

\begin{definition}
An integer $Q\geq 1$ \emph{closes} the T-matrix of a WZW model
if $Q\,T_j$ is an odd multiple of $\tfrac{1}{4}$ for every primary~$j$
(equivalently, $\cos(2\pi Q T_j)=0$ for all~$j$).
\end{definition}

\begin{theorem}[T-matrix closure at $(E_6)_1$ level]
\label{P15:thm:QgroupE6}
The minimum closing integer for $(E_6)_1$ using~\eqref{P15:eq:TE6} is
$Q_{\mathrm{group}}^{(3)}=3$.
\end{theorem}
\begin{proof}
Primary $\mathbf{27}$: $T=\tfrac{5}{12}$.
$Q\cdot\tfrac{5}{12}=(2n+1)/4$ requires $Q\cdot 5=3(2n+1)$,
so $Q\in\{3,9,15,\ldots\}$.
Primary $\mathbf{1}$: $T=-\tfrac{1}{4}$.
$Q\cdot(-\tfrac{1}{4})=(2m+1)/4$ requires $Q$ odd:
$Q\in\{1,3,5,\ldots\}$.
The minimum in $\{3,9,\ldots\}\cap\{1,3,5,\ldots\}$ is $Q=3$.
Verification: $3\cdot(-\tfrac{1}{4})=-\tfrac{3}{4}$ (odd/4);\;
$3\cdot\tfrac{5}{12}=\tfrac{15}{12}=\tfrac{5}{4}$ (odd/4).
\qed
\end{proof}

\begin{theorem}[T-matrix closure at SU(3)$_1$ level]
\label{P15:thm:Qgroup3}
The minimum closing integer for $\SU(3)_1$ using~\eqref{P15:eq:TSU3} is
also $Q_{\mathrm{group}}^{(3)}=3$.
\end{theorem}
\begin{proof}
Primary $(1,0)$: $T=\tfrac{1}{4}$, requires $Q$ odd.
Primary $(0,0)$: $T=-\tfrac{1}{12}$, requires $Q\in\{3,9,15,\ldots\}$.
Minimum: $Q=3$.
Verification: $3\cdot(-\tfrac{1}{12})=-\tfrac{1}{4}$ (odd/4);\;
$3\cdot\tfrac{1}{4}=\tfrac{3}{4}$ (odd/4).
\qed
\end{proof}

\begin{remark}[Robustness of $Q_{\mathrm{group}}^{(3)}=3$]
The closing integer $Q=3$ is obtained independently at both levels
of the WZW chain~\eqref{P15:eq:confemb}.
It equals $|2T|/8=3$ (where~8 is the order of the
octahedral faces).
The agreement between the $(E_6)_1$ and $\SU(3)_1$ results
confirms the internal consistency of the WZW chain.
\end{remark}

%══════════════════════════════════════════════════════════════════════════════
\section{Both Quark Charges from the E$_6$ Structure}
\label{P15:sec:quarkcharge}
%══════════════════════════════════════════════════════════════════════════════

\begin{theorem}[Both quark charges from the WZW chain]
\label{P15:thm:charge}
Both fractional quark charges emerge at different levels of the
WZW chain~\eqref{P15:eq:confemb}, with no free parameters:
\begin{enumerate}[noitemsep,label=(\roman*)]
\item $(E_6)_1$ level: $h_{\mathbf{27}}=\tfrac{2}{3}$
  gives the up-quark charge $|Q_u|=\tfrac{2}{3}$.
\item $\SU(3)_1$ level ($\SU(3)_L$ electroweak factor):
  \begin{equation}
    h_{(1,0)} = \frac{p^2+pq+q^2+3p+3q}{12}\bigg|_{p=1,q=0}
              = \frac{4}{12} = \frac{1}{3}
    \label{P15:eq:hdown}
  \end{equation}
  gives the down-quark charge $|Q_d|=\tfrac{1}{3}$.
\item These are consistent via branching~\eqref{P15:eq:27branch}:
  $h_{\mathbf{3}}+h_{\bar{\mathbf{3}}}
  =\tfrac{1}{3}+\tfrac{1}{3}=\tfrac{2}{3}=h_{\mathbf{27}}$.
\end{enumerate}
\end{theorem}

\begin{proof}
(i) Equation~\eqref{P15:eq:h27}.
(ii) Substitution into~\eqref{P15:eq:hdown}.
(iii) The conformal embedding~\eqref{P15:eq:confemb} and
branching~\eqref{P15:eq:27branch}: the $(\mathbf{3},\bar{\mathbf{3}},\mathbf{1})$
component has weight $h_{(1,0)}+h_{(0,1)}=2/3=h_{\mathbf{27}}$. \qed
\end{proof}

\begin{remark}[Both charges, no free parameters]
Neither $|Q_u|=2/3$ nor $|Q_d|=1/3$ is an input; both are outputs
of the McKay correspondence and conformal embedding.
The ratio $|Q_u|/|Q_d|=2$ follows from the $E_6\supset\SU(3)^3$
branching.
\end{remark}

\subsection{Electric charge from the trinification generator}
\label{P15:sec:charge_generator}

The identifications $|Q_d|=h_{(1,0)}$ and $|Q_u|=h_{\mathbf{27}}$
in Theorem~\ref{P15:thm:charge} equate a conformal weight (an $L_0$
eigenvalue) with an electric charge (a $U(1)$ representation
label) -- a priori different types of quantity. We now show this
equality is forced by the trinification embedding itself, with no
further identification.

\begin{theorem}[Electric charge from the trinification generator]
\label{P15:thm:charge_generator}
Let $\SU(3)_L\subset(E_6)_1$ be the electroweak factor of the
trinification $E_6\supset\SU(3)_C\times\SU(3)_L\times\SU(3)_R$, and
let
\begin{equation}
  Q = T_{3L}+\frac{1}{\sqrt3}T_{8L}
  \label{P15:eq:Qgen}
\end{equation}
be the unique electric-charge generator of\/ $\SU(3)_L$:
the linear combination $\alpha T_{3L}+\beta T_{8L}$ of Cartan
generators that is traceless over the fundamental $\mathbf{3}_L$
and compatible with the $\mathbb{Z}_3$ centre of\/ $\SU(3)_L$
(charges are integer multiples of $\tfrac{1}{3}$).
These two requirements uniquely force $\alpha=1$,
$\beta=1/\!\sqrt{3}$: the eigenvalue of $T_{8L}$ on the
third weight state is $-1/\!\sqrt{3}$, so
$\beta\cdot(-1/\!\sqrt{3})=-\tfrac{1}{3}$ gives $\beta=1/\!\sqrt{3}$
with no free parameter, and then $\alpha=1$ follows from
tracelessness.
The formula $Q=T_{3L}+T_{8L}/\!\sqrt{3}$ is therefore an output
of the $\SU(3)_L$ group theory of the trinification embedding,
not an external input; its agreement with the Standard Model
hypercharge assignment is a consequence, not a postulate.
On the three weight states of the fundamental $\mathbf{3}_L$,
\begin{equation}
  (T_{3L},T_{8L}) \in
  \Bigl\{\bigl(\tfrac12,\tfrac{1}{2\sqrt3}\bigr),\;
         \bigl(-\tfrac12,\tfrac{1}{2\sqrt3}\bigr),\;
         \bigl(0,-\tfrac{1}{\sqrt3}\bigr)\Bigr\},
  \label{P15:eq:T3T8}
\end{equation}
giving
\begin{equation}
  Q \in \Bigl\{\tfrac{2}{3},\,-\tfrac{1}{3},\,-\tfrac{1}{3}\Bigr\},
  \label{P15:eq:trinification_charges}
\end{equation}
the standard $(u,d,D)$ charge assignment of one trinification
family. In the Cartan inner product normalised so that
$(\alpha,\alpha)=2$ for roots of $\SU(3)_L$ (the level-1 WZW
convention used in~\eqref{P15:eq:hdown}), the charge generator's
coefficient vector is
\begin{equation}
  \mu = 2\,\omega_1,
  \label{P15:eq:mu_omega1}
\end{equation}
where $\omega_1$ is the highest weight of $\mathbf{3}_L$ -- i.e.\
the electric-charge generator is (twice) the highest weight of the
very representation whose $\SU(3)_1$ conformal weight is
$h_{(1,0)}$.
\end{theorem}

\begin{proof}
Direct evaluation of~\eqref{P15:eq:Qgen} on~\eqref{P15:eq:T3T8} gives
\eqref{P15:eq:trinification_charges}.
Writing $(T_{3L},T_{8L})$ as Cartan coordinates, the highest-weight
state of $\mathbf{3}_L$ has $\omega_1=(\tfrac12,\tfrac{1}{2\sqrt3})$,
and the coefficient vector of $Q$ in~\eqref{P15:eq:Qgen} is
$\mu=(1,\tfrac{1}{\sqrt3})=2\omega_1$. \qed
\end{proof}

\begin{corollary}[Quark charges equal conformal weights, derived]
\label{P15:cor:charge_eq_weight}
Since $h_{(1,0)}=\tfrac12(\omega_1,\omega_1)=\tfrac13$
(level-1 WZW conformal-weight formula) and
$Q(\omega_1)=(\omega_1,\mu)=2(\omega_1,\omega_1)=\tfrac23$
by~\eqref{P15:eq:mu_omega1}, both quark charges and both conformal
weights of Theorem~\ref{P15:thm:charge} are computed from the single
number $(\omega_1,\omega_1)_{A_2}=\tfrac23$:
\begin{equation}
  \boxed{|Q_d| = h_{(1,0)} = \tfrac13,
  \qquad
  |Q_u| = Q(\omega_1) = 2h_{(1,0)} = h_{\mathbf{27}} = \tfrac23.}
  \label{P15:eq:charge_eq_weight}
\end{equation}
No identification beyond Theorem~\ref{P15:thm:charge_generator} is
required: the equality of conformal weight and electric charge is
a consequence of the trinification embedding together with the
$A_2$ representation theory of Section~\ref{P15:sec:wzw}, and holds
only because $N=3$ (Remark~\ref{P15:rem:N3_unique}).
\end{corollary}

\begin{remark}[Why $N=3$ is essential]
\label{P15:rem:N3_unique}
For $\SU(N)_1$, $h_{\mathrm{fund}}=(N-1)/(2N)$, while the
$\mathbb{Z}_N$-centre-completing $U(1)$ charge of the fundamental
-- the minimal charge consistent with
$U(N)=(\SU(N)\times U(1))/\mathbb{Z}_N$ -- is $1/N$. These coincide,
\begin{equation}
  h_{\mathrm{fund}} = \frac{1}{N}
  \quad\Longleftrightarrow\quad
  \frac{N-1}{2N}=\frac{1}{N}
  \quad\Longleftrightarrow\quad
  N=3.
  \label{P15:eq:N3unique}
\end{equation}
The identification of Corollary~\ref{P15:cor:charge_eq_weight} is
therefore not a generic feature of $\SU(N)_1$ WZW models -- it is
special to $N=3$, the value independently forced by the McKay
correspondence $\twoT\leftrightarrow\widehat{E}_6$
(Section~\ref{P15:sec:wzw}). The $\QH=3$ sector and the
``conformal weight $=$ electric charge'' coincidence are
simultaneously $N=3$ phenomena.
\end{remark}

\begin{remark}[Leading-order mass degeneracy]
In $\SU(3)_1$: $T_{(1,0)}=T_{(0,1)}=\tfrac{1}{4}$ (identical
T-matrix phases); by the same logic as Papers~III--IV~\cite{Paper3,Paper4},
this predicts $m_u=m_d$ at leading order.
In $(E_6)_1$: $T_{\mathbf{27}}=T_{\overline{\mathbf{27}}}=\tfrac{5}{12}$
(same degeneracy).
The prediction $m_u\approx m_d$ is consistent with the measured
ratio $m_u/m_d\approx0.46$, the smallest quark mass ratio,
generated by subleading isospin-breaking effects.
\end{remark}

%══════════════════════════════════════════════════════════════════════════════
\section{Virial Fixed Point and Energy Formula}
\label{P15:sec:virial}
%══════════════════════════════════════════════════════════════════════════════

\begin{theorem}[Universal virial identity]
\label{P15:thm:virial}
At the density-feedback virial saddle of any sector $n$,
$V^{*\,(n)}\equiv J_{\mathrm{fb}}^{(n)}/\Ja^{(n)}=\vp$,
and consequently
\begin{equation}
  \Kfb^{(n)} = 2\vp\,\Ja^{(n)},
  \qquad
  \Efb^{(n)} = 2N_n^2/\vp^{17}.
  \label{P15:eq:virial}
\end{equation}
\end{theorem}

\begin{proof}
For the $\QH=3$ sector: numerical evaluation of the density-feedback
functional at the initial profile gives
$V^{(3)}(0)\approx 2.58>\vp=1.618$,
while $V^{(3)}(\infty)=0<\vp$ trivially.
By the intermediate value theorem applied to the continuous,
monotone function $V^{(3)}(b)$ (Paper~III~\cite{Paper3}
Theorem~3.5), there exists a unique $b^*$ with $V^{(3)}(b^*)=\vp$.
This is the same IVT argument as for $\QH=2$; the feedback kernel
depends only on local field magnitudes, not on the topological sector.

With $V^*=\vp$ and $\ms=3-\vp$ (exact):
\begin{align*}
  \Kfb &= \Ja(1+\ms\vp) = \Ja(1+(3-\vp)\vp) \\
       &= \Ja(1+3\vp-\vp^2) = \Ja(1+3\vp-\vp-1) = 2\vp\,\Ja
\end{align*}
using $\vp^2=\vp+1$.

For the energy: if $\vp^9\sqrt{\Ja\Jfour}=N_n$ and $\Jfour=r_n\Ja$,
then $\Ja\Jfour=r_n\Ja^2=N_n^2/\vp^{18}$, so the virial ratio
$r_n$ cancels and
$\Efb = \Kfb\cdot\Jfour = 2\vp\Ja\cdot r_n\Ja
= 2\vp\cdot N_n^2/(\vp^{18})=2N_n^2/\vp^{17}$.
\qed
\end{proof}

\begin{remark}[Sector independence of the energy formula]
Equation~\eqref{P15:eq:virial} shows that $\Efb^{(n)}=2N_n^2/\vp^{17}$
depends only on the integer $N_n$, not on the virial ratio $r_n$.
The sector geometry changes $r_n$ but not the overall energy scale
(at fixed $N_n$).
\end{remark}

%══════════════════════════════════════════════════════════════════════════════
\section{Trefoil Tube Geometry: Frenet--Serret Transformation}
\label{P15:sec:geometry}
%══════════════════════════════════════════════════════════════════════════════

\subsection{The Q$_H$=2 profile as the Q$_H$=3 tube profile}

The BS profile $f(\rho)=2\arctan(\rho^{-C^*})$ is the same function
for both $\QH=2$ (torus) and $\QH=3$ (trefoil tube).
The one-dimensional profile integrals
\begin{equation}
  I_{2a} = \tfrac{32}{15},\qquad
  I_{4\mathrm{tube}} = \frac{8}{5C^{*2}},\qquad
  I_{2\mathrm{iso}} = \text{(same form)}
  \label{P15:eq:1Dints}
\end{equation}
depend only on the tube cross-section and are independent of the
centre-curve geometry.
The Frenet--Serret transformation
\begin{equation}
  \mathbf{x}(\rho,\chi,s) = \Gamma(s) + \rho\bigl[\mathbf{n}(s)\cos\chi
    + \mathbf{b}(s)\sin\chi\bigr]
  \label{P15:eq:FS}
\end{equation}
maps the torus parametrisation exactly onto a tube of radius $\rho$
around any smooth curve $\Gamma(s)$, where $\mathbf{n}(s)$,
$\mathbf{b}(s)$ are the Frenet--Serret normal and binormal.
Replacing $\Gamma$ with the trefoil centre curve $T_{2,3}$
gives the \emph{exact} $\QH=3$ tube embedding.

\subsection{Jacobian and leading-order energy integrals}

The volume element under~\eqref{P15:eq:FS} is
\begin{equation}
  dV = (1 - \rho\,\kappa(s)\cos\chi)\,\rho\,d\rho\,d\chi\,ds,
  \label{P15:eq:Jacobian}
\end{equation}
where $\kappa(s)$ is the curvature of $\Gamma(s)$.
This is the \emph{same Jacobian} computed for the thick-torus
deformation in Paper~XIV~\cite{Paper14}, with $1/R_0$ replaced
by the variable curvature $\kappa(s)$ of the trefoil.

For the standard embedding of $T_{2,3}$ with $R_0=3$, $r_0=0.874$:
\begin{equation}
  \langle\kappa\rangle = 0.335\approx 1/R_0 = 0.333,
  \qquad
  \langle\kappa^2\rangle = 0.120\approx 1/R_0^2 = 0.111.
  \label{P15:eq:kap_stats}
\end{equation}
The mean-square curvature of the trefoil is nearly identical to
that of the torus circle at $R_0=3$.
The arc length is $L_3=41.26$, compared to $L_2=2\pi R_0=18.85$,
giving $L_3/L_2=2.19$.

\begin{remark}[Physical origin of $r_0$: the satellite construction]
\label{P15:rem:r0_origin}
Unlike $R_0$, which is fixed by the Y-junction string-length
condition $3R_0=N_3^2=9$ (Proposition~\ref{P15:prop:yjunction}), the
minor radius $r_0$ is \emph{not} an independent geometric input to
the $\QH=3$ construction.  Paper~XVIII~\cite{Paper18}
(Proposition~3.5, %\ref{P18:prop:satellite} \ref{P18:rem:satellite_physical}
Remark~3.6) identifies $T_{2,3}$ as the satellite knot obtained by threading a
$\QH=1$ unknot through the solid-torus neighbourhood $N(\gamma_1)$ of
one component of the $\QH=2$ Hopf link, with the $(2,3)$-cable
framing.  The trefoil's tube is therefore built \emph{around} the
already-existing $\QH=2$ tube, not as an independent $\QH=3$ energy
minimum: $r_0$ should inherit its scale from $N(\gamma_1)$, i.e.\
from the $\QH=2$ sector's own profile geometry, rather than from a
$\QH=3$-internal saddle condition.  This is consistent with the
near-coincidence
\begin{equation}
  r_0 \;\approx\; \frac{R_0}{C_2^*} \;=\; \frac{3}{3.4318} \;=\; 0.874172,
  \label{P15:eq:r0_satellite}
\end{equation}
matching the construction value to $0.020\%$, where $C_2^*=3.4318$
is the $\QH=2$ sector's own modified-BPS profile-sharpness constant
(Paper~I~\cite{Paper1}, Theorem~3.1): the unique real root of
$C(2C+1)/[(C^2+1)(3C-1)]=2^{4/3}/\vp^5$.  If $1/C_2^*$ is identified
with the dimensionless tube radius of $N(\gamma_1)$ in units of
$R_0$, equation~\eqref{P15:eq:r0_satellite} converts the otherwise
unexplained input $r_0=0.874$ into a quantity inherited from an
already-derived $\QH=2$ constant.  This identification is
\emph{motivated} by the satellite construction but not yet
\emph{derived} from it: Paper~XVIII~\cite{Paper18} states that the cable framing is
``dictated by'' the $\QH=2$ BPS condition without giving the explicit
map from $N(\gamma_1)$'s radius to $r_0$.  Supplying that map is an
open problem; until it is closed, $r_0=0.874$ should be read as
consistent with, but not proved by, equation~\eqref{P15:eq:r0_satellite}.
\end{remark}

\subsection{Torsion vanishes at leading order}

\begin{theorem}[Exact torsion cancellation for symmetric profiles]
\label{P15:thm:torsion_cancels}
In the Frenet--Serret frame~\eqref{P15:eq:FS}, the torsion $\tau(s)$
acts as a rigid rotation $\chi\to\chi-\int_0^s\tau(s')\,ds'$ of the
tube cross-section coordinate.
For any profile $f=f_0(\rho)$ depending only on $\rho$
(i.e.\ $\partial f/\partial\chi=0$), this rotation
is invisible at \emph{every order}: the torsion makes
zero contribution to $J_4^{(3)}/J_{2a}^{(3)}$ at all orders in $\rho\tau$.
\end{theorem}

\begin{proof}
The torsion introduces the substitution $\chi\to\chi+\phi(s)$ where
$\phi(s)=\int_0^s\tau\,ds'$.  For any function $g(\rho,\chi)=h(\rho)$
independent of $\chi$ (which is the case for the BS profile $f_0(\rho)$
and all integrands built from it),
$\int_0^{2\pi}h(\rho)\,d(\chi+\phi(s)) = \int_0^{2\pi}h(\rho)\,d\chi$.
The torsion rotation is therefore invisible at \emph{every} order,
not merely at first order.
In contrast, the Frenet normal $\mathbf{n}$ itself changes direction
with torsion, but since the Hopf curvature components
$F_{\rho\chi}=f_0'(\rho)\sin(f_0)\cos(f_0)/\rho$
involve only $\partial f/\partial\rho$ (not $\partial f/\partial\chi$
or $\partial f/\partial s$), they are independent of $\tau$.
\qed
\end{proof}

\begin{remark}[Scope of the cancellation: leading order only]
\label{P15:rem:torsion_scope}
The cancellation proved above is exact for the leading-order,
$\chi$-independent profile $f_0(\rho)$ alone: torsion is a pure
rotation of $\chi$, and a $\chi$-independent integrand is invariant
under any such rotation, to all orders in $\rho\tau$, with no
further hypotheses needed. This is the content of
Theorem~\ref{P15:thm:torsion_cancels} and it is correct as stated.

It does \emph{not} follow that torsion is absent from the full
trefoil energy at every order. Once the curvature-sourced first
correction $f_1(\rho)\cos\chi$ is included --- the same correction
responsible for the $\langle\kappa^2\rangle$ term already present
in~\eqref{P15:eq:r3formula} --- torsion has a non-trivial $\chi$-dependent
function to act on. In the Frenet--Serret tube metric, torsion enters
as the off-diagonal term $g_{\chi s}=\rho^2\tau$, exact and absent
only because the torus has $\tau\equiv 0$. The induced metric
determinant $\sqrt{g}=\rho(1-\rho\kappa\cos\chi)$ is exactly
torsion-independent (a rotation cannot change a volume element), so
the Jacobian~\eqref{P15:eq:Jacobian} and the leading profile equation for
$f_0$ are unaffected. But the inverse metric is not, and torsion acts
non-trivially once $f_1(\rho)\cos\chi\neq 0$: the leading surviving
contribution to the $\chi$-averaged gradient energy is
\begin{equation}
  \bigl\langle\delta|\nabla f|^2\bigr\rangle_\chi
  = \tfrac{1}{2}\,f_1(\rho)^2\,\kappa(s)^2\,\tau(s)^2\,\rho^2
  + O(\rho^3),
  \label{P15:eq:torsion_correction}
\end{equation}
proportional to $\kappa^2\tau^2$, not $\tau^2$ alone: torsion is
gated by curvature and vanishes whenever either does. Along the
trefoil centre curve, $\langle\kappa^2\tau^2\rangle\approx 0.0053$,
roughly $4\%$ of $\langle\kappa^2\rangle=0.120$ already present
in~\eqref{P15:eq:r3formula}. The full boundary-value problem for the
curvature- and torsion-sourced correction, including a structural
finding that the relevant indicial exponents near $\rho=0$ are
complex, is developed in Paper~XVI. The cancellation of
Theorem~\ref{P15:thm:torsion_cancels} should therefore be read as a
leading-order statement only, sound by exactly the mechanism it
proves, and not as evidence that torsion is irrelevant to the
trefoil's full $O(\rho^2)$ energy.
\end{remark}

\begin{corollary}[Leading-order trefoil formula]
At leading order in $\rho/R_0$,
\begin{equation}
  \boxed{\frac{J_4^{(3)}}{J_{2a}^{(3)}}
    = \frac{3}{4C_3^{*2}} + \frac{\langle\kappa^2\rangle}{2}
    + O\!\left(\frac{\rho}{R_0}\right)^{\!2},}
  \label{P15:eq:r3formula}
\end{equation}
where $\langle\kappa^2\rangle=0.120$ is the arc-length-weighted
mean-square curvature of the trefoil~\eqref{P15:eq:kap_stats},
replacing $1/(2R_0^2)$ in the $\QH=2$ formula.
Torsion makes no contribution at this order
(Theorem~\ref{P15:thm:torsion_cancels}; see Remark~\ref{P15:rem:torsion_scope}
for its $O(\rho^2\kappa^2\tau^2)$ contribution at next order).
\end{corollary}

\begin{corollary}[Bogomolny parameter from C$_3^*$]
\label{P15:cor:lambda3}
Combining Theorem~\ref{P15:thm:virial} ($\Kfb=2\vp\,\Ja$)
with~\eqref{P15:eq:r3formula} ($J_4^{(3)}/\Ja^{(3)}=r_3$), the
leading-order Bogomolny parameter of the $\QH=3$ sector is
\begin{equation}
  \boxed{\lam_3
  = \frac{\Kfb}{J_4^{(3)}}
  = \frac{2\vp}{r_3}
  = \frac{2\vp}{\dfrac{3}{4C_3^{*2}}+\dfrac{\langle\kappa^2\rangle}{2}},}
  \label{P15:eq:lambda3}
\end{equation}
at the same order as~\eqref{P15:eq:r3formula}.
Equivalently, the saddle parameter $C_3^*$ is determined by
$\lam_3$ alone:
\begin{equation}
  C_3^* = \sqrt{\frac{3}{4\!\left(\dfrac{2\vp}{\lam_3}-\dfrac{\langle\kappa^2\rangle}{2}\right)}}.
  \label{P15:eq:C3star}
\end{equation}
If the Bogomolny parameter is sector-independent,
$\lam_3=\vp^6=\lam_2$, then $r_3=2/\vp^5=0.1803$ and
\begin{equation}
  C_3^* = 2.497.
  \label{P15:eq:C3leading}
\end{equation}
\end{corollary}

\begin{proof}
Equation~\eqref{P15:eq:lambda3} follows immediately from
$\Kfb/J_4^{(3)} = 2\vp\Ja^{(3)}/(r_3\Ja^{(3)}) = 2\vp/r_3$
with $r_3$ given by~\eqref{P15:eq:r3formula}.
Equation~\eqref{P15:eq:C3star} is the algebraic inverse.
For~\eqref{P15:eq:C3leading}: set $\lam_3=\vp^6$ so $2\vp/\vp^6=2\vp^{-5}$;
then $3/(4C_3^{*2})=2/\vp^5-0.060=0.1203$, giving
$C_3^{*2}=3/(4\times0.1203)=6.234$ and $C_3^*=2.497$ at $O(\langle\kappa^2\rangle)$.
Including the $O(\langle\kappa^4\rangle)$ Beta-function correction
(Proposition~\ref{P15:prop:kappa4}): $3/(4C_3^{*2})=0.1201$,
giving $C_3^*=2.5062$ (shift of $0.392\%$).
\qed
\end{proof}



\begin{remark}[Tube separation]
Numerical computation of the trefoil $T_{2,3}$ gives minimum
inter-tube distance $d_{\min}=1.748$, while the tube
cross-section radius is $1/C^*_2\approx0.291$.
The ratio $d_{\min}\cdot C^*_2\approx 6.0$ confirms that the
tubes are well-separated and the thin-tube approximation is valid
at leading order.
\end{remark}

%══════════════════════════════════════════════════════════════════════════════
\section{Crossing Regions, Inner-Wall Analogy, and Colour}
\label{P15:sec:crossings}
%══════════════════════════════════════════════════════════════════════════════

\subsection{The trefoil inner wall}

For the $\QH=2$ torus, the density-feedback Shafranov shift
concentrates the condensate toward the inner wall ($r<R_0$),
producing an inner/outer density ratio of $3.6\times$ at $\rho=1$
(Paper~XIV~\cite{Paper14}).
The mechanism is curvature-driven: higher $\kappa$ on the inner wall
increases the energy density, and feedback concentrates the
condensate there.

For the $\QH=3$ trefoil, the same mechanism operates but the
geometry is different.
At the three crossing regions of $T_{2,3}$, the curvature reaches
its maximum $\kappa_{\max}=0.399$ (compared to the mean $\langle\kappa\rangle=0.335$).
By Theorem~\ref{P15:thm:torsion_cancels}, torsion does not contribute
at this order.
The Shafranov shift~\eqref{P15:eq:Jacobian} is therefore enhanced
at the crossing regions:
\begin{equation}
  \Delta^{(3)}(\rho,s) > \Delta^{(3)}(\rho,s')
  \quad\text{whenever}\quad \kappa(s) > \kappa(s').
  \label{P15:eq:crossing_shift}
\end{equation}
The three crossing regions carry \emph{enhanced condensate density},
playing exactly the role of the inner wall for the $\QH=2$ torus.

\begin{proposition}[Crossing regions as inner wall]
\label{P15:prop:crossings}
The three crossing regions of the trefoil $T_{2,3}$ are the
$\QH=3$ analogue of the torus inner wall:
\begin{enumerate}[noitemsep,label=(\roman*)]
\item They are the loci of maximum curvature and hence maximum
  Shafranov shift along the trefoil.
\item They carry enhanced condensate density relative to the
  far sections of the tube, by the same curvature-driven mechanism
  as the $3.6\times$ inner/outer ratio of $\QH=2$.
\item They are related by the $\mathbb{Z}_3$ rotational symmetry
  of the trefoil --- a symmetry absent in $\QH=2$.
\end{enumerate}
\end{proposition}

\subsection{Three concentrations as three colours}

The $\mathbb{Z}_3$ symmetry of the crossing concentrations provides
the geometric identification of quark colours.
The three condensate concentration peaks are:
\begin{itemize}[noitemsep]
\item Related by $120^\circ$ rotation (the $\mathbb{Z}_3$ outer
  automorphism of $\SU(3)_1$).
\item Individually distinguishable by their position along the
  trefoil --- they label ``colour~1'', ``colour~2'', ``colour~3''.
\item Non-interchangeable by any continuous deformation that
  preserves the trefoil knot type.
\end{itemize}
This is the condensate imprint of colour charge: each quark
colour corresponds to one crossing-region concentration of the
$\QH=3$ Hopfion.

\begin{remark}[Geometric centre is vacuum]
The geometric centre of the trefoil at $R_0=3$ is the origin,
at minimum distance $2.13$ from the nearest tube point --- more
than $7\times$ the tube radius $1/C^*\approx 0.29$.
The origin lies in the field vacuum.
The interior of the trefoil (the Seifert solid) has a condensate
\emph{minimum} at its geometric centre and maxima at the three
crossing regions --- the inverse of an incorrect intuition that
the ``centre'' would be filled.
\end{remark}

\begin{remark}[Three concentrations: exact count]
\label{P15:rem:3or6}
Each of the three crossing regions of $T_{2,3}$ consists of two strands
(over and under), raising the question of whether the energy concentration
count is $3$ or $6$ (or some larger multiple of $3$).
A direct calculation settles this.
The crossings occur at parameter values $t_1=\pi/6+2n\pi/3$ (where $\cos 3t=0$),
with partner strand at $t_2=t_1+\pi$.
The curvature at any crossing satisfies
\begin{equation}
  \kappa(t_1)=\kappa(t_1+\pi)=0.3991
  \label{P15:eq:kappa_equal}
\end{equation}
for all three crossings simultaneously.
This is not a numerical coincidence: the map
$t\mapsto t+\pi$ combined with $(x,y,z)\mapsto(x,y,-z)$
is an exact isometry of $T_{2,3}$ that exchanges the over- and under-strands
at every crossing, forcing $\kappa_{\mathrm{over}}=\kappa_{\mathrm{under}}$.
(The same map proves $|\Gamma'(t)|^2 = 9r_0^2\sin^2 3t + 4R_0^2$,
independent of the sign of $\sin 3t$, so the two strands also have equal
arc-length speed.)

Since over and under strands carry identical curvature, they contribute
equal Shafranov-shifted energy density.
Each crossing region is a single energy concentration with $\mathbb{Z}_2$-symmetric
over/under contributions.
The count is therefore exactly $\mathbf{3}$: one per crossing, as claimed in
Proposition~\ref{P15:prop:crossings}.
Higher multiples ($6, 9, 12, \ldots$) require either broken over/under symmetry
(which is absent by the above isometry) or additional crossing regions
(the single minimum $d_{\min}=2r_0$ rules out more than three).

\smallskip
\noindent\textit{Additional geometric facts verified analytically:}
$d_{\min}=|\Gamma(t)-\Gamma(t+\pi)|=2r_0=1.748$ for all $t$ (constant!);
$\kappa_{\max}=0.399$ at the crossing regions; $\langle\kappa^2\rangle=0.120$;
$\langle\kappa\rangle=0.335$; $L_3=41.26$.
\end{remark}

%══════════════════════════════════════════════════════════════════════════════
\section{Colour Confinement from Profile Topology}
\label{P15:sec:confinement}
%══════════════════════════════════════════════════════════════════════════════

\begin{proposition}[Confinement from boundary conditions]
\label{P15:prop:confinement}
Colour confinement is already encoded in the boundary conditions
of the Paper~I~\cite{Paper1} profile:
\begin{equation}
  f(0) = \pi,\qquad f(\infty) = 0.
  \label{P15:eq:BC}
\end{equation}
An isolated open flux strand (one tube of the trefoil with free
endpoints) requires $f$ to interpolate $0\to\pi\to 0$ along its
length, with the field returning to the vacuum at both ends.
The gradient energy of such a configuration diverges linearly
with strand length.
Only closed configurations ($f(\text{end})=f(\text{start})$)
have finite energy, requiring the three trefoil tubes to link
into the closed knot $T_{2,3}$.
\end{proposition}

\begin{proof}
For an open strand of length $L$, the tube must transition
$f:0\to\pi$ at the source end and $f:\pi\to 0$ at the sink end.
The transition cost per unit length in the bulk is bounded below
by the BPS gradient energy $\sim C^*/L$ (from the profile integrals).
As $L\to\infty$, the total energy $\sim C^*$ per crossing is finite,
but the contribution from maintaining $f=\pi$ in the bulk
scales as $L$ through the $(f')^2$ term in $J_{2a}$.
The energy therefore diverges as $L\to\infty$, prohibiting
isolated open strands.
\qed
\end{proof}

\begin{remark}[Q$_H$=2 vs Q$_H$=3 confinement]
The confinement mechanism is the same in both sectors:
the profile boundary conditions force closed tubes.
For $\QH=2$ (one tube, closed torus): the tube closes on itself;
this is the topologically protected soliton of Paper~I.
For $\QH=3$ (three tubes, trefoil): the three tubes must link to
close; an isolated colour charge cannot form a finite-energy
closed configuration alone.
Only colour-neutral combinations --- three tubes linked in
$T_{2,3}$ (baryons) or one tube plus its anti-tube (mesons in
$\QH=3\oplus\QH=-3$) --- are energetically allowed.
The quantitative statement (linear potential coefficient) requires
knowledge of $C^*_3$ from the $\QH=3$ saddle.
\end{remark}

\begin{remark}[Absence of a unique Bogomolny parameter and confinement]
\label{P15:rem:no_lambda3}
For the $\QH=2$ torus, the single inner-wall energy concentration
defines a unique saddle with Bogomolny parameter $\lam_2=\vp^6$.
For the $\QH=3$ trefoil, the three crossing-region concentrations
(Proposition~\ref{P15:prop:crossings}, Remark~\ref{P15:rem:3or6}) produce a
\emph{three-lobed energy landscape} in configuration space, with no
single saddle.
Extensive numerical gradient flow on $80^3$ grids with multiple
algorithms --- direct $\Kfb$ minimisation, fixed-$\lam$ flows at
$\lam=\vp^6$ and $\lam=\vp^5$, and projected-gradient flow at
fixed $J_4$ --- consistently shows $\Kfb/J_4$ traversing $\vp^6$
without converging there.
The trajectory of $\Kfb/J_4$ forms a $V$-shape under $\Kfb$
minimisation (minimum $\approx 3.55$, crossing $\vp^6$ twice), and
an inverted-$U$ shape under fixed-$\lam$ flows (peak values $\approx
21\vp^6$ and $12\vp^5$ for the two test values).
Neither flow converges to a unique saddle at $\Kfb/J_4=\vp^k$ for
any integer $k$.

This behaviour is not a numerical failure.
It reflects the physical structure of the three-lobed landscape:
separating one tube of the trefoil requires simultaneously traversing
all three crossing-region energy barriers, which costs more energy
than the field can supply.
The $\QH=3$ condensate is therefore not a stable isolated soliton
admitting a unique Bogomolny parameter, but a \emph{transient
collective field state} whose three tubes must co-exist.
This is confinement seen in the energy landscape: the Bogomolny
parameter $\lam_2=\vp^6$ is definable for the electron because it
is a stable isolated soliton; the analogous $\lam_3$ is not uniquely
definable for the quark sector because quarks are never isolated.

The appropriate characterisation of the $\QH=3$ sector is through
the \emph{global} invariants $N_3=3$ (T-matrix closure,
Theorem~\ref{P15:thm:Qgroup3}) and $\Kfb=2\vp\Ja$ (virial,
Theorem~\ref{P15:thm:virial}), rather than through a local saddle value.
Numerical comparison of the fixed-$\lam$ flows establishes the
bound $\lam_3>\vp^5=11.09$: at equal integration times, the flow
with $\lam=\vp^5$ preserves $J_4$ better than $\lam=\vp^6$, and
higher $\lam$ causes faster $J_4$ collapse, setting a lower bound
on any effective sector Bogomolny parameter.
The Corollary~\ref{P15:cor:lambda3} formula $C_3^*=\sqrt{3/[4(2\vp/\lam_3
-\langle\kappa^2\rangle/2)]}$ remains the best analytical prediction:
since $\lam_3=\vp^6$ is proved (Theorem~\ref{P15:thm:sector_assignment}),
this gives $C_3^*=2.5062$ at $O(\langle\kappa^4\rangle)$
(Proposition~\ref{P15:prop:kappa4}), which
awaits confirmation from a topology-preserving saddle-finding
algorithm (open problem~(iii) in Section~\ref{P15:sec:open}).
\end{remark}

\begin{proposition}[Y-junction geometry: the Fermat--Steiner energy balance]
\label{P15:prop:yjunction}
The three crossing midpoints
$\mathbf{m}_k = \tfrac{1}{2}[\Gamma(t_k)+\Gamma(t_k+\pi)]$,
$k=1,2,3$, satisfy
\begin{equation}
  |\mathbf{m}_k| = R_0,\qquad
  |\mathbf{m}_k - \mathbf{m}_j| = R_0\sqrt{3}\quad(j\neq k),\qquad
  \sum_{k=1}^3 \mathbf{m}_k = \mathbf{0}.
  \label{P15:eq:yjunction}
\end{equation}
The three midpoints lie at the vertices of a regular equilateral
triangle in the $z=0$ plane, all at distance $R_0$ from the origin.
The origin is their centroid and the Fermat--Steiner point (the
unique point minimising total distance to the three vertices) of
this equilateral triangle, with all three arms of equal length~$R_0$.
\end{proposition}

\begin{proof}
At the crossing $t_k$ (where $\cos 3t_k=0$), we showed in the proof
of Remark~\ref{P15:rem:3or6} that $\Gamma(t_k)=(R_0\cos 2t_k,\, R_0\sin 2t_k,\, r_0)$
and $\Gamma(t_k+\pi)=(R_0\cos 2t_k,\, R_0\sin 2t_k,\, -r_0)$.
Hence $\mathbf{m}_k=(R_0\cos 2t_k,\, R_0\sin 2t_k,\, 0)$,
so $|\mathbf{m}_k|=R_0$.
With $2t_1=\pi/3$, $2t_2=5\pi/3$, $2t_3=3\pi$ the three unit vectors
are at $60^\circ, 300^\circ, 180^\circ$ respectively, summing to zero.
The side length $|\mathbf{m}_1-\mathbf{m}_2|=R_0\,|\mathrm{e}^{i\pi/3}
-\mathrm{e}^{-i\pi/3}|=R_0\sqrt{3}$.
For an equilateral triangle the Fermat--Steiner point equals the
centroid, which is $\mathbf{0}$.
\qed
\end{proof}

\begin{remark}[The Y-junction as the missing energy term]
\label{P15:rem:yjunction_energy}
Proposition~\ref{P15:prop:yjunction} identifies the origin as the
\emph{Y-junction} (Fermat--Steiner vertex) of the three-quark baryon.
The full baryon energy has three contributions:
\begin{equation}
  E_{\mathrm{baryon}} = \underbrace{K_{\mathrm{fb}}\times J_4}_{\text{profile}}
  + \underbrace{\sigma\,\cdot\,3R_0}_{\text{string tension}}
  + \underbrace{E_\cup}_{\text{junction}},
  \label{P15:eq:baryon_energy}
\end{equation}
where $\sigma$ is the colour string tension and $3R_0$ is the
total Y-junction string length (three equal arms of length $R_0$
meeting at $120^\circ$ at the origin).

All numerical flows described in Remark~\ref{P15:rem:no_lambda3}
minimised only the profile term $K_{\mathrm{fb}}\times J_4$.
This explains why $K_{\mathrm{fb}}/J_4$ never converged: each
crossing region experiences an additional inward force
$\sigma\,\hat{\mathbf{r}}_k$ (unit vector toward origin) from
the string tension, which modifies the Bogomolny condition to
\begin{equation}
  J_4\,\frac{\delta K_{\mathrm{fb}}}{\delta\mathbf{n}}
  + K_{\mathrm{fb}}\,\frac{\delta J_4}{\delta\mathbf{n}}
  = -\sigma\,\hat{\mathbf{r}}_k\,\frac{\delta J_4}{\delta\mathbf{n}}.
  \label{P15:eq:modified_bogomolny}
\end{equation}
Without the string term the three crossing-region saddles compete,
producing the observed oscillating $K_{\mathrm{fb}}/J_4$ trajectories.
With the string term they are locked at $|\mathbf{m}_k|=R_0$ and
the full baryon energy~\eqref{P15:eq:baryon_energy} has a proper saddle.

By contrast, the $\QH=2$ torus has no string network (one tube,
no Y-junction), so its energy is purely the profile term
$K_{\mathrm{fb}}\times J_4$, which has a unique saddle at $\lam_2=\vp^6$.

The total Y-junction string length $3R_0$ carries a further
significance: $3R_0/R_0 = N_3 = 3$.
The normalisation identity $\vp^9\sqrt{\Ja^{(3)}J_4^{(3)}}=3$
(Theorem~\ref{P15:thm:norm3}) thus encodes not only the T-matrix closure
integer but also the Y-junction arm count and total string length.
\end{remark}

%══════════════════════════════════════════════════════════════════════════════
\section{The Subgroup $2T\subset 2I$: The 24-Cell, Colour, and Generation Structure}
\label{P15:sec:2T2I}
%══════════════════════════════════════════════════════════════════════════════

\subsection{The subgroup relationship}

\begin{theorem}[$2T\subset 2I$]
\label{P15:thm:subgroup}
The binary tetrahedral group $\twoT$ (order~24) is a subgroup of
the binary icosahedral group $\twoI$ (order~120).
\end{theorem}

\begin{proof}
The icosahedral group $I\cong A_5$ (order~60) contains the
tetrahedral group $T\cong A_4$ (order~12) as a subgroup,
$A_4\subset A_5$, since every even permutation of four elements
is an even permutation of five.
Lifting via the 2:1 covering $\SU(2)\to\mathrm{SO}(3)$ gives
$\twoT = 2A_4\subset 2A_5 = \twoI$.
\qed
\end{proof}

\begin{remark}[Index and structure]
The index $[\twoI:\twoT]=120/24=5$ equals the pentagonal
symmetry number $k+2\big|_{k=3}=5$ of the WZW level.
The five cosets of $\twoT$ in $\twoI$
correspond to the five lepton primaries after including the
vacuum $j=0$, and the restriction $\twoI\to\twoT$ is the algebraic
content of passing from the lepton sector ($\QH=2$) to the
quark sector ($\QH=3$).
\end{remark}

\subsection{The 24-cell: geometry of the quark sector}

The 120 elements of $\twoI$ are the vertices of the 600-cell in $\mathbb{R}^4$,
whose edge-to-circumradius ratio $1/\vp$ is the geometric origin of the
golden ratio throughout the lepton sector.
An exact parallel holds for $\twoT$:

\begin{proposition}[24-cell structure of $\twoT$]
\label{P15:prop:24cell}
The 24 elements of $\twoT\subset\SU(2)$, as unit quaternions,
are the vertices of the \emph{24-cell} in $\mathbb{R}^4$.
They decompose as two sets:
\begin{align}
  \text{Set A (quaternion subgroup }Q_8\text{):}&\quad
    \{\pm 1, \pm\mathbf{i}, \pm\mathbf{j}, \pm\mathbf{k}\},
  \label{P15:eq:Q8_vertices}\\
  \text{Set B (Hurwitz units):}&\quad
    \bigl\{\tfrac12(\pm1\pm\mathbf{i}\pm\mathbf{j}\pm\mathbf{k})\bigr\}
    \quad\text{(all 16 sign combinations).}
  \label{P15:eq:Hurwitz_vertices}
\end{align}
The 24-cell has edge length~$1$ and circumradius~$1$; its
\emph{edge-to-circumradius ratio is exactly~$1$},
compared to $1/\vp$ for the 600-cell.
The distance spectrum from any vertex is
$\{1\,(8\text{ neighbours}),\;\sqrt{2}\,(6),\;\sqrt{3}\,(8),\;2\,(1\text{ antipodal})\}$.
\end{proposition}

\begin{proof}
All 24 elements satisfy $|q|^2=1$ (circumradius~1).
Direct computation: $|1-\mathbf{i}|=\sqrt{2}$ (Set~A nearest)
but $|1-\tfrac12(1+\mathbf{i}+\mathbf{j}+\mathbf{k})|=1$ (between Set~A and Set~B),
confirming edge length~1.
The full distance spectrum follows by symmetry and explicit enumeration.
\qed
\end{proof}

\begin{remark}[The $F_4$ root system and the Coxeter number for quarks]
\label{P15:rem:F4}
The 24 short roots of the exceptional Lie algebra $F_4$ form exactly the
24-cell. The analogue of the $E_8$ McKay correspondence for $\twoI$ is:
\begin{equation}
  \twoI \;\leftrightarrow\; E_8:\quad
  |\twoI|=120 = \tfrac12\cdot|R(E_8)|=\tfrac12\cdot 240,
  \qquad
  \twoT \;\leftrightarrow\; F_4:\quad
  |\twoT|=24 = \tfrac12\cdot|R(F_4)|=\tfrac12\cdot 48.
  \label{P15:eq:2T_F4}
\end{equation}
The relevant Coxeter numbers are $h(E_8)=30$ (entering the lepton Yukawa
$y_e=e^{-25\pi/6}$ of Paper~IV~\cite{Paper4}) and $h(F_4)=12$ for quarks.
The ratio $h(F_4)/h(E_8)=12/30=2/5=(k+2)^{-1}$ at $k=3$ is consistent
with the index $[\twoI:\twoT]=5=k+2$.
The dual Coxeter number $h^\vee(F_4)=9=N_3^2$ matches the $\QH=3$
normalisation integer squared.
\end{remark}

\begin{remark}[Integer normalisation from the 24-cell]
The 600-cell's edge-to-circumradius ratio $1/\vp$ produces the
golden-ratio arithmetic of the $\QH=2$ sector: $\vp^9\sqrt{\Ja J_4}=10$.
The 24-cell's ratio of exactly~$1$ (no $\vp$) corresponds to the
$\QH=3$ normalisation identity $\vp^9\sqrt{\Ja^{(3)}J_4^{(3)}}=3$
having an \emph{integer} right-hand side.
The $\vp$ factors in the $\QH=3$ identity come entirely from the
condensate (the $\vp^9$ prefactor), not from the quark-sector group.
This reflects the 24-cell's rational arithmetic: there is no
golden ratio in the geometry of $\twoT$.
Since the colour gauge structure $\SU(3)_1$ arises from the conformal
embedding $(E_6)_1\supset\SU(3)_1^{\times3}$ (Section~\ref{P15:sec:embed}),
and $(E_6)_1$ is itself the McKay dual of the $\vp$-free group $\twoT$,
the colour algebra inherits no $\vp$-dependence at the group-theoretic
level; this is consistent with $\vp$ being wholly absent from QCD's
own gauge structure (the Gell-Mann algebra, the running of $\alpha_s$,
and the colour-confinement mechanism), even though $\vp$ enters the
condensate energy scale $\Efb^{(3)}=2N_3^2/\vp^{17}$ that sets the
absolute mass scale (open problem~(iv) of Section~\ref{P15:sec:open}).
\end{remark}

\subsection{Conjugacy classes and character table}

The 7 conjugacy classes of $\twoT$, with class sizes summing to 24,
are listed in Table~\ref{P15:tab:2T_classes}.

\begin{table}[h!]
\centering
\renewcommand{\arraystretch}{1.25}
\begin{tabular}{clcc}
\toprule
Class & Representative & Size & Order \\
\midrule
$C_1$ & $+1$ & 1 & 1 \\
$C_2$ & $-1$ & 1 & 2 \\
$C_3$ & $\mathbf{i}$ (and $\pm\mathbf{j},\,\pm\mathbf{k},\,-\mathbf{i}$) & 6 & 4 \\
$C_4$ & $\tfrac12(1+\mathbf{i}+\mathbf{j}+\mathbf{k})$ & 4 & 6 \\
$C_5$ & $\tfrac12(1+\mathbf{i}+\mathbf{j}-\mathbf{k})$ & 4 & 6 \\
$C_6$ & $\tfrac12(-1+\mathbf{i}+\mathbf{j}+\mathbf{k})$ & 4 & 3 \\
$C_7$ & $\tfrac12(-1+\mathbf{i}+\mathbf{j}-\mathbf{k})$ & 4 & 3 \\
\bottomrule
\end{tabular}
\caption{Conjugacy classes of $\twoT$.
Classes $C_1$ and $C_2$ form the centre $Z(\twoT)=\{\pm1\}\cong\mathbb{Z}_2$.
Classes $C_4,C_5$ (order~6) and $C_6,C_7$ (order~3) are the Hurwitz units.}
\label{P15:tab:2T_classes}
\end{table}

Let $\omega=e^{2\pi i/3}$.
The full character table of $\twoT$ is given in Table~\ref{P15:tab:2T_chars}.

\begin{table}[h!]
\centering
\renewcommand{\arraystretch}{1.25}
\setlength{\tabcolsep}{6pt}
\begin{tabular}{lccccccc}
\toprule
Irrep & $C_1$ & $C_2$ & $C_3$ & $C_4$ & $C_5$ & $C_6$ & $C_7$ \\
\midrule
$\Gamma_1$ (dim~1) & $1$ & $1$ & $1$ & $1$ & $1$ & $1$ & $1$ \\
$\Gamma_\omega$ (dim~1) & $1$ & $1$ & $1$ & $\omega$ & $\omega^2$ & $\omega^2$ & $\omega$ \\
$\Gamma_{\omega^2}$ (dim~1) & $1$ & $1$ & $1$ & $\omega^2$ & $\omega$ & $\omega$ & $\omega^2$ \\
$\Gamma_2$ (dim~2) & $2$ & $-2$ & $0$ & $1$ & $1$ & $-1$ & $-1$ \\
$\Gamma_{2\omega}$ (dim~2) & $2$ & $-2$ & $0$ & $\omega$ & $\omega$ & $-\omega^2$ & $-\omega^2$ \\
$\Gamma_{2\omega^2}$ (dim~2) & $2$ & $-2$ & $0$ & $\omega^2$ & $\omega^2$ & $-\omega$ & $-\omega$ \\
$\Gamma_3$ (dim~3) & $3$ & $3$ & $-1$ & $0$ & $0$ & $0$ & $0$ \\
\bottomrule
\end{tabular}
\caption{Character table of $\twoT$, $\omega=e^{2\pi i/3}$.
The three \emph{fermionic} irreps $\Gamma_2,\Gamma_{2\omega},\Gamma_{2\omega^2}$
(those with $\chi(C_2)=-\dim$, i.e.\ that change sign under $q\mapsto -q$)
are the three quark colour representations.
Sum of squares: $3\cdot1+3\cdot4+9=24=|\twoT|$. \checkmark}
\label{P15:tab:2T_chars}
\end{table}

\subsection{Branching rules $2I\to 2T$}

Using the inner product $n_\rho = |\twoT|^{-1}\sum_c\chi_j(C_c)\,\overline{\chi_\rho(C_c)}\,|C_c|$
with the $\mathrm{SU}(2)$ characters
$\chi_j(C_c)=(\pm1)^{2j}\sin((2j+1)\theta_c/2)/\sin(\theta_c/2)$
(where $\theta_c$ is the rotation angle of class $C_c$ and the sign
is $(-1)^{2j}$ when $\mathrm{Re}(q)<0$),
the complete branching rules are:

\begin{proposition}[Branching rules $2I\to 2T$, proved by character theory]
\label{P15:prop:branching}
\begin{equation}
\renewcommand{\arraystretch}{1.3}
\begin{array}{rll}
  j=0\phantom{/2}& \to & \Gamma_1 \\[2pt]
  j=\tfrac12 & \to & \Gamma_2 \\[2pt]
  j=1\phantom{/2} & \to & \Gamma_3 \\[2pt]
  j=\tfrac32 & \to & \Gamma_{2\omega}\oplus\Gamma_{2\omega^2} \\[2pt]
  j=2\phantom{/2} & \to & \Gamma_\omega\oplus\Gamma_{\omega^2}\oplus\Gamma_3 \\[2pt]
  j=\tfrac52 & \to & \Gamma_2\oplus\Gamma_{2\omega}\oplus\Gamma_{2\omega^2}
\end{array}
\label{P15:eq:branching_table}
\end{equation}
Dimension check: $1,2,3,4,5,6$ match the $\mathrm{SU}(2)$ dimensions $2j+1$.
All multiplicities are $0$ or $1$: the restriction $2I\to 2T$ is
multiplicity-free for $j\le\tfrac52$.
\end{proposition}

\begin{proof}
Evaluate $\chi_j$ on each class representative (Table~\ref{P15:tab:2T_classes})
using the $\mathrm{SU}(2)$ character formula, then compute $n_\rho$ by the inner
product over the 7 classes with sizes $(1,1,6,4,4,4,4)$.
All multiplicities evaluate to integers~0 or~1; the dimension sums
match by construction. Full numerical verification in the companion
script \texttt{paper15\_simple\_current.py}.
\qed
\end{proof}

\begin{remark}[Key physical features of the branching rules]
Three entries of~\eqref{P15:eq:branching_table} have direct physical content.

\emph{(i) $j=\tfrac12\to\Gamma_2$:} the $\QH=1$ fundamental maps to
colour~1. This grounds the bosonization identification (Paper~III
Proposition~10.6) across both sectors.

\emph{(ii) $j=\tfrac32\to\Gamma_{2\omega}\oplus\Gamma_{2\omega^2}$:}
the simple current of $\mathrm{SU}(2)_3$ decomposes into
colours~2 and~3 only, not colour~1.
The three colours together appear only in the triple fusion
$j=\tfrac12\otimes j=\tfrac12\otimes j=\tfrac12$
(which contains $j=\tfrac52\to\Gamma_2\oplus\Gamma_{2\omega}\oplus\Gamma_{2\omega^2}$,
all three colours simultaneously), confirming the three-tube structure of the trefoil.

\emph{(iii) $j=0,\tfrac12,1\to\Gamma_1,\Gamma_2,\Gamma_3$:}
the three lepton-generation levels of $\mathrm{SU}(2)_3$
each map to a \emph{distinct} irrep of $\twoT$.
For quarks, the generation label $(j_g,\Gamma_{2\omega^c})$ carries
T-matrix phase $T_{j_g}^{\mathrm{SU}(2)_3}\in\{-\tfrac{3}{40},\tfrac{3}{40},\tfrac{13}{40}\}$
with $Q_3=3$, giving inter-generation T-phase ratios that are $O(1)$ corrections;
the large mass hierarchy comes from the $\vp$-tower levels (open problem~(iv)(c)).
\end{remark}

\subsection{Representation theory and physical assignments}

\begin{proposition}[Three colours from three fermionic irreps]
\label{P15:prop:colours}
The three $2$-dimensional fermionic irreps
$\Gamma_2,\Gamma_{2\omega},\Gamma_{2\omega^2}$ of $\twoT$,
related by the $\mathbb{Z}_3$ outer automorphism of $E_6$,
are the three quark colours.
They arise from restricting the spinor representations of $\twoI$
to the subgroup $\twoT$:
\begin{equation}
  \twoI(\text{spinor}, \dim=2)\;\big\downarrow_{\twoT}\;
  \longrightarrow\; \Gamma_{2\omega^k},
  \qquad k\in\{0,1,2\},
  \label{P15:eq:branching}
\end{equation}
with index $k$ labelling the quark colour.
\end{proposition}

The branching rule~\eqref{P15:eq:branching} provides the algebraic
bridge: quark colours are the $\mathbb{Z}_3$-twisted restrictions
of the lepton-sector spinor irrep of $\twoI$.
The $\mathbb{Z}_3$ automorphism that permutes the three colours
is precisely the outer automorphism of $E_6$ that permutes the
three legs of the $\widehat{E}_6$ McKay diagram.

\subsection{Quark charge and leading-order mass degeneracy}

\begin{theorem}[Charge and mass from $\SU(3)_1$]
\label{P15:thm:massdegen}
Within the $\SU(3)_1$ WZW framework for the $\QH=3$ sector:
\begin{enumerate}[noitemsep,label=(\roman*)]
\item The fractional electric charge of the quark is
  $h_{(1,0)}=\tfrac{1}{3}$~(Theorem~\ref{P15:thm:charge}).
\item The T-matrix exponents satisfy
  $T_{(1,0)}=T_{(0,1)}=\tfrac{1}{4}$: the fundamental
  $\mathbf{3}$ and anti-fundamental $\bar{\mathbf{3}}$ have
  identical T-matrix phases.
\item At leading order in the WZW framework, this degeneracy
  implies $m_u=m_d$ (up--down quark mass equality).
\end{enumerate}
\end{theorem}

\begin{proof}
(i) Theorem~\ref{P15:thm:charge}.
(ii) Direct substitution into~\eqref{P15:eq:TSU3}.
(iii) In Papers~III--IV~\cite{Paper3,Paper4} distinct T-phases give
distinct masses; identical phases $T_{(1,0)}=T_{(0,1)}$ give
identical leading-order masses. \qed
\end{proof}

\begin{remark}[Consistency with observation]
The prediction $m_u\approx m_d$ at leading order is consistent
with the measured ratio $m_u/m_d\approx 0.46$ (the smallest among six quarks).
The large hierarchies $m_c/m_u\sim10^3$, $m_t/m_u\sim10^5$
involve the $\vp$-tower levels and the quark Yukawa coupling,
deferred to Paper~XVI~\cite{Paper16} (open problem~(iv)(c) of
Section~\ref{P15:sec:open}).
\end{remark}

%══════════════════════════════════════════════════════════════════════════════
\section{The Normalisation Identity}
\label{P15:sec:norm}
%══════════════════════════════════════════════════════════════════════════════

\begin{theorem}[Normalisation identity]
\label{P15:thm:norm3}
Given the virial identity $\Kfb=2\vp\Ja$ (Theorem~\ref{P15:thm:virial}),
the sector energy formula $\Efb^{(3)}=2N_3^2/\vp^{17}$, and the
thermodynamic axiom $\Efb^{(3)}/V_3=\rCMB=10$ (Paper~XII~\cite{Paper12}):
if $\lam_3=\vp^6$, then
\begin{equation}
  \boxed{\vp^9\sqrt{\Ja^{(3)}\Jfour^{(3)}}
  = 3 = Q_{\mathrm{group}}^{(3)}.}
  \label{P15:eq:norm3}
\end{equation}
\end{theorem}

\begin{proof}
$\lam_3=\vp^6 \Rightarrow r_3=2/\vp^5$
(Corollary~\ref{P15:cor:lambda3}).
$\Efb=(4/\vp^4)\Ja^2$
(from $\Kfb=2\vp\Ja$, $\Jfour=r_3\Ja$).
$\Efb=2N_3^2/\vp^{17} \Rightarrow \Ja=N_3/(\sqrt{2}\,\vp^{13/2})$.
Then
$\vp^9\sqrt{\Ja\Jfour}=\vp^9\cdot\Ja\cdot\sqrt{r_3}
=\vp^9\cdot N_3/(\sqrt{2}\,\vp^{13/2})\cdot\sqrt{2}/\vp^{5/2}
=N_3=3$.
\qed
\end{proof}

\begin{corollary}[Condensate volume as derived quantity]
\label{P15:cor:V3}
Under the same hypothesis $\lam_3=\vp^6$,
the trefoil condensate volume is
\begin{equation}
  V_3 = \frac{N_3^2}{5\vp^{17}} = \frac{9}{5\vp^{17}}.
  \label{P15:eq:V3}
\end{equation}
This is a \emph{derived consequence} of $\lam_3=\vp^6$,
not an independent assumption.
\end{corollary}

\begin{proof}
$V_3=\Efb/\rCMB=(2N_3^2/\vp^{17})/10=N_3^2/(5\vp^{17})=9/(5\vp^{17})$.
\qed
\end{proof}

\begin{remark}[Why $N_3=3$ is forced]
The CMB factor $\rCMB=10\Lcond^4$ is sector-independent.
The integer $N_n=Q_{\mathrm{group}}^{(n)}$ is determined independently
by the T-matrix closure (Theorem~\ref{P15:thm:Qgroup3}).
The normalisation identity and volume formula are therefore
consequences of the condensate physics once $\lam_3=\vp^6$ is proved;
they do not need to be assumed.
\end{remark}

\begin{theorem}[Logical equivalence of the two closure conditions]
\label{P15:thm:equivalence}
The following two statements are logically equivalent,
given the thermodynamic axiom and energy formula:
\begin{enumerate}[noitemsep,label=(\Roman*)]
\item $\lam_3 = \vp^6$ (sector-independent Bogomolny parameter).
\item $\vp^9\sqrt{\Ja^{(3)}\Jfour^{(3)}} = 3$ (normalisation identity).
\end{enumerate}
Either implies, as a corollary (Corollary~\ref{P15:cor:V3}),
the volume formula $V_3=9/(5\vp^{17})$.
\end{theorem}

\begin{proof}
\textbf{(I)$\Rightarrow$(II)} is Theorem~\ref{P15:thm:norm3} (four-line proof above).

\textbf{(II)$\Rightarrow$(I).}
If $\vp^9\sqrt{\Ja\Jfour}=3$, then $N_3=3$.
From $\Efb=(4/\vp^4)\Ja^2=2N_3^2/\vp^{17}$:
$\Ja=3/(\sqrt{2}\,\vp^{13/2})$ and $\Jfour=9/(\vp^{18}\Ja)=3\sqrt{2}/\vp^{23/2}$.
Hence $r_3=\Jfour/\Ja=2/\vp^5$, and
$\lam_3=2\vp/r_3=\vp^6$. $\checkmark$

\smallskip\noindent
(Numerical check: $\vp^9\sqrt{\Ja\Jfour}=3.00000000$ to eight figures.)
\qed
\end{proof}

\begin{remark}[Volume formula as geometric self-consistency check]
The trefoil condensate volume satisfies the geometric identity
\begin{equation}
  V_3 = L_3\cdot\pi C_3^{*-2}\cdot\mathcal{F}(\kappa,\tau,\rho/R_0)
  \label{P15:eq:V3formula}
\end{equation}
where $L_3=41.26$ and $\mathcal{F}$ accounts for curvature,
torsion (negligible at this order, $O(\kappa^2\tau^2)$,
Remark~\ref{P15:rem:torsion_scope}), and finite-tube effects.
With $C_3^*=2.5062$ (Proposition~\ref{P15:prop:kappa4}),
equation~\eqref{P15:eq:V3formula} should equal $9/(5\vp^{17})$
(Corollary~\ref{P15:cor:V3}), providing a geometric self-consistency check.
\end{remark}

\begin{proposition}[$O(\langle\kappa^4\rangle)$ correction via Beta functions]
\label{P15:prop:kappa4}
In normalised tube coordinates $\hat\rho = \rho C_3^*$ the BS profile is
$f_0(\hat\rho)=2\arctan(\hat\rho^{-1})$, so that
$\sin f_0 = 2\hat\rho/(1+\hat\rho^2)$ and
$f_0'(\hat\rho)=-2/(1+\hat\rho^2)$.
Setting $t=\hat\rho^2$, the profile integrals reduce to Beta functions:
\begin{equation}
  I_n \;\equiv\;
  \int_0^\infty \sin^{2n}\!f_0\,(f_0')^2\,\hat\rho\,d\hat\rho
  \;=\; 2^{2n+1}B\!\left(n+1,\,n+1\right).
  \label{P15:eq:In_beta}
\end{equation}
In particular $I_2 = \tfrac{16}{15}$ (the normalised $J_{2a}$ integral).

The $O(\rho^4\kappa^4)$ term in the Jacobian expansion of $J_4^{(3)}$,
after averaging $\langle\cos^4\chi\rangle = \tfrac{3}{8}$ over the tube
angle, involves the moment
\begin{equation}
  I^{(4)} \;\equiv\;
  \int_0^\infty \sin^4\!f_0\,(f_0')^2\,\hat\rho^4\,\hat\rho\,d\hat\rho
  \;=\; 32\,B(5,1) \;=\; \frac{32}{5},
  \label{P15:eq:Ik4}
\end{equation}
giving the correction
\begin{equation}
  \delta r_3^{(\kappa^4)}
  \;=\; \frac{3}{8}\,\langle\kappa^4\rangle\,\frac{1}{C_3^{*4}}\,
        \frac{I^{(4)}}{I_2}
  \;=\; \frac{9\,\langle\kappa^4\rangle}{4\,C_3^{*4}}.
  \label{P15:eq:dr_kappa4}
\end{equation}
For $T_{2,3}$ with $\langle\kappa^4\rangle=0.0162$ and $C_3^*=2.497$:
\begin{equation}
  \delta r_3^{(\kappa^4)}
  \;=\; \frac{9\times0.0162}{4\times2.497^4}
  \;=\; 0.000938
  \;=\; 0.520\%\text{ of }r_3.
  \label{P15:eq:dr_kappa4_num}
\end{equation}
Since $\lam_3=\vp^6$ is proved independently and $r_3=2/\vp^5$ is fixed,
the correction shifts $C_3^*$ (not $\lam_3$) via
$3/(4C_3^{*2}) = r_3 - \langle\kappa^2\rangle/2 - \delta r_3^{(\kappa^4)}$:
\begin{equation}
  \boxed{C_3^* = 2.5062 \quad\bigl(O(\langle\kappa^4\rangle)\text{ corrected}\bigr),}
  \label{P15:eq:C3star_kappa4}
\end{equation}
a shift of $0.392\%$ from the $O(\langle\kappa^2\rangle)$ value $2.4965$.
\end{proposition}

\begin{proof}
The Beta-function identity $\int_0^\infty t^{a-1}/(1+t)^{a+b}\,dt = B(a,b)$
with $t=\hat\rho^2$ converts each profile integral to a standard Beta function.
For $I_n$: writing the integrand $\sin^{2n}\!f_0\,(f_0')^2\,\hat\rho
=4^{n+1}\hat\rho^{2n+1}/(1+\hat\rho^2)^{2n+2}$ and substituting
$t=\hat\rho^2$ ($d\hat\rho=dt/(2\sqrt t)$) gives
$I_n=2^{2n+1}\int_0^\infty t^{n}/(1+t)^{2n+2}\,dt$, so $a=n+1$, $b=n+1$,
giving equation~\eqref{P15:eq:In_beta};
$I_2 = 2^5 B(3,3) = 32\times\tfrac{1}{30}=\tfrac{16}{15}$, using
$B(n+1,n+1)=\Gamma(n+1)^2/\Gamma(2n+2)$ at $n=2$.
For $I^{(4)}$: the $\hat\rho^4$ power shifts $a$ by $2$, giving $a=5,b=1$ and
$B(5,1)=\Gamma(5)\Gamma(1)/\Gamma(6)=24\cdot1/120=1/5$; hence $I^{(4)}=32/5$.
The ratio $I^{(4)}/I_2 = (32/5)/(16/15) = 6$; inserting into~\eqref{P15:eq:dr_kappa4}
and simplifying ($\tfrac38\times6=\tfrac94$) gives~\eqref{P15:eq:dr_kappa4}.
\qed
\end{proof}

%══════════════════════════════════════════════════════════════════════════════
\section{Sector Hierarchy}
\label{P15:sec:discussion}
%══════════════════════════════════════════════════════════════════════════════

\subsection{The lepton--quark hierarchy}

The normalisation integers encode the sector hierarchy:
\begin{equation}
  N_2 = Q_{\mathrm{group}}^{(2)} = 10 \;\leftrightarrow\; \text{leptons},
  \qquad
  N_3 = Q_{\mathrm{group}}^{(3)} = 3 \;\leftrightarrow\; \text{quarks}.
\end{equation}
The energy ratio $\Efb^{(3)}/\Efb^{(2)} = N_3^2/N_2^2 = 9/100$
(from Theorem~\ref{P15:thm:virial}) relates the hadronic and leptonic
condensate energy scales without free parameters.

%══════════════════════════════════════════════════════════════════════════════
\section{Condensate Universality and the $\lam_3=\vp^6$ Conjecture}
\label{P15:sec:lambda3_conjecture}
%══════════════════════════════════════════════════════════════════════════════

By Theorem~\ref{P15:thm:equivalence}, proving any one of the three
equivalent conditions (I)--(III) closes all the others
simultaneously.  The most tractable entry point is
$\lam_3=\vp^6$ via the condensate universality argument below.

\subsection{Why $\lam_3=\vp^6$ is expected}

\begin{proposition}[Condensate universality of $\lam_3$]
\label{P15:prop:lambda3}
If the $\QH=3$ density-feedback functional $\Efb^{(3)}$ admits
a topology-preserving saddle with virial fixed point $V^*=\vp$
(proved in Theorem~\ref{P15:thm:virial}), and if the $2D$ EL saddle
profile satisfies $r_3 = \Jfour^{(3)}/\Ja^{(3)} = 2/\vp^5$
at the saddle, then $\lam_3=\vp^6$.
\end{proposition}

\begin{proof}
From Theorem~\ref{P15:thm:virial}: $\Kfb=2\vp\Ja$.
Thus $\lam_3 = \Kfb/\Jfour = 2\vp\Ja/\Jfour = 2\vp/r_3$.
Setting $r_3 = 2/\vp^5$ gives $\lam_3 = 2\vp/(2/\vp^5)=\vp^6$.
\qed
\end{proof}

\noindent
The open question therefore reduces to a single geometric statement:
\begin{equation}
  \frac{\Jfour^{(3)}}{\Ja^{(3)}}\bigg|_{\mathrm{EL\;saddle}} = \frac{2}{\vp^5}.
  \label{P15:eq:r3_conjecture}
\end{equation}
This is the $\QH=3$ analogue of Paper~III's Linking Scale Conjecture,
whose proof occupied Theorems~4.3, 4.4, and~10.1 of that paper.

\subsection{The condensate-universality argument}

The proof of the $\QH=2$ Linking Scale Conjecture (Paper~III,
Theorems~4.3--4.4 and~10.1) had three steps:
\begin{enumerate}[noitemsep,label=(\alph*)]
\item \emph{$V^*=\vp$ (IVT):} proved sector-independently.
  For $\QH=3$: \textbf{done} (Theorem~\ref{P15:thm:virial}).
\item \emph{$\Kfb=2\vp\Ja$ (algebra):} follows from $\vp^2=\vp+1$.
  For $\QH=3$: \textbf{done} (Theorem~\ref{P15:thm:virial}).
\item \emph{$r_2=J_4^{(2)}/J_{2a}^{(2)}=2^{4/3}/\vp^5$ at the 2D EL saddle:}
  proved in Paper~III by showing that the additional Hopf curvature
  components $F_{12},F_{23}$ (contributing $\approx30\%$ of $J_4$
  in the converged torus profile) cancel exactly in the ratio
  $J_4/J_{2a}$, via the $\SU(2)_3$ $S$-matrix self-duality
  $S_{\frac12,\frac12}=S_{0,0}$ at $k=3$
  (Paper~II Proposition~7.1 and its proof).
\end{enumerate}

For $\QH=3$, steps~(a) and~(b) are done.
Step~(c) requires the trefoil analogue: show that all
crossing-region curvature contributions to $J_4^{(3)}$ cancel
in the ratio $r_3$, enforced by the same $\SU(2)_3$ condensate
$S$-matrix self-duality.

\begin{remark}[Why the same self-duality applies]
The condensate background is $\SU(2)_3$ for all sectors.
The self-duality $S_{\frac12,\frac12}=S_{0,0}$
(Paper~II Proposition~7.1) is a property of the condensate WZW model,
not of the individual soliton sector.
For $\QH=2$ it controlled the $F_{12},F_{23}$ components of the
torus; for $\QH=3$ it should control the analogous crossing-region
integrals of the trefoil.
The Rellich--Kondrachov compactness argument that extends
$V^*=\vp$ to finite $R_0$ (Paper~III Theorem~9.4) applies unchanged
to the trefoil tube geometry: the tube is a compact
manifold-with-boundary, and the $\QH=3$ charge constraint prevents
mass escape at infinity exactly as for $\QH=2$.
\end{remark}

\begin{remark}[Relation to the thin-tube prediction]
The leading-order formula~\eqref{P15:eq:r3formula} gives
$r_3 = 3/(4C_3^{*2})+\langle\kappa^2\rangle/2$.
Setting $r_3=2/\vp^5$ yields $C_3^*=2.497$
(Corollary~\ref{P15:cor:lambda3} with $\lam_3=\vp^6$).
By Theorem~\ref{P15:thm:torsion_cancels}, torsion makes zero
contribution at this order (Remark~\ref{P15:rem:torsion_scope}).
The $O(\langle\kappa^4\rangle)$ correction shifts $C_3^*$ to
$2.5062$ ($0.392\%$ shift; Proposition~\ref{P15:prop:kappa4}) and
leaves $\lam_3$ self-consistent.
The statement~\eqref{P15:eq:r3_conjecture} is consistent with all
analytical results established so far.
\end{remark}

\subsection{Status}

Proposition~\ref{P15:prop:lambda3} reduces the problem
to the geometric identity~\eqref{P15:eq:r3_conjecture}.
Once proved (e.g.\ by identifying the crossing-region integrals
with specific $\SU(2)_3$ $S$-matrix elements via the Paper~III
methodology, and applying Rellich--Kondrachov to the trefoil tube),
the proof closes $\lam_3=\vp^6$ and, as derived consequences
(Corollary~\ref{P15:cor:V3} and Theorem~\ref{P15:thm:equivalence}),
$V_3=9/(5\vp^{17})$ and $\vp^9\sqrt{\Ja^{(3)}\Jfour^{(3)}}=3$
simultaneously.

\subsection{The geometric expansion and its WZW interpretation}
\label{P15:sec:geom_expansion}

We now lay out the proof strategy for~\eqref{P15:eq:r3_conjecture}
at the level of a structured argument, identifying each step
that is proved versus each step that requires the Paper~III
methodology applied to the trefoil.

\subsubsection*{Structure of the exact $r_3$ formula}

The Frenet--Serret volume element $(1-\rho\kappa\cos\chi)\rho\,d\rho\,d\chi\,ds$
acts differently on $J_4^{(3)}$ and $J_{2a}^{(3)}$:

\begin{proposition}[Exact structure of $r_3$ corrections]
\label{P15:prop:r3_exact}
Let $f_0(\rho)$ be any axially-symmetric profile.
\begin{enumerate}[label=\textup{(\roman*)}]
\item $J_{2a}^{(3)}$ receives \emph{no} curvature correction from the
  Jacobian: $\delta J_{2a}^{\kappa} = 0$ exactly, since
  $\int_0^{2\pi}\cos\chi\,d\chi=0$.
\item The leading curvature correction to $J_4^{(3)}$ comes from
  the $s$-derivative Hopf component $F_{s\chi}\sim \kappa\cos\chi \cdot \sin^2\!f_0$,
  whose squared modulus chi-averages to
  $\langle F_{s\chi}^2\rangle_\chi = \tfrac12\kappa^2 \cdot \sin^4\!f_0 \cdot (f_0')^2$.
\item The ratio therefore has the exact expansion
  \begin{equation}
    r_3 = \frac{3}{4C_3^{*2}}
    + \frac{\langle\kappa^2\rangle_s}{2}\cdot\Gamma(C_3^*)
    + O(\langle\kappa^4\rangle_s),
    \label{P15:eq:r3_exact_expansion}
  \end{equation}
  where $\Gamma(C)$ is a dimensionless profile-integral ratio
  (explicitly computable via Beta functions for the BS profile family).
\end{enumerate}
\end{proposition}

\begin{proof}
Part~(i): $J_{2a} = \int \sin^4\!f\,[{(f')}^2+\sin^2\!f/\rho^2]\cdot
(1-\rho\kappa\cos\chi)\,\rho\,d\rho\,d\chi\,ds$.
The integrand is independent of $\chi$; the $\cos\chi$ factor in the
Jacobian integrates to zero over $[0,2\pi]$.

Part~(ii): In the FS frame with $f=f_0(\rho)$, the non-zero
$s$-derivative component arises from the rotation of the normal
$\mathbf{n}$ at rate $\kappa(s)$:
$F_{s\chi} = \kappa(s)\cos\chi \cdot \sin^2\!f_0 \cdot f_0'(\rho)$.
Squaring and averaging: $\langle F_{s\chi}^2\rangle_\chi = \tfrac12\kappa^2\sin^4\!f_0(f_0')^2$.
Since $\langle\cos^2\chi\rangle=\tfrac12$ and the integrand has no
further $\chi$-dependence.

Part~(iii) follows by combining (i) and (ii) in the ratio $J_4/J_{2a}$
and expanding in powers of $\langle\kappa^{2n}\rangle_s$.
\qed
\end{proof}

\subsubsection*{The algebraic target and the fusion-count identity}

A key structural observation closes the argument at the level of
the WZW self-consistency:

\begin{lemma}[Fusion-count ratio of sector targets]
\label{P15:lem:fusion_ratio}
The target ratios for the two sectors satisfy
\begin{equation}
  \frac{r_3^*}{r_2^*}
  = \frac{2/\vp^5}{2^{4/3}/\vp^5}
  = \frac{1}{2^{1/3}}
  = N_{\mathrm{fus}}^{-1/k},
  \label{P15:eq:fusion_ratio}
\end{equation}
where $N_{\mathrm{fus}}=2$ (the fusion count; see Paper~II Proposition~9.2)
and $k=3$ (the WZW level).
Equivalently, $r_3^* = r_2^*/N_{\mathrm{fus}}^{1/k}$.
\end{lemma}

\begin{proof}
Direct substitution: $r_2^* = 2^{4/3}/\vp^5$ (Paper~III) and
$r_3^* = 2/\vp^5$ (from $\lam_3=\vp^6$ via Proposition~\ref{P15:prop:lambda3}).
The ratio is $2/2^{4/3} = 2^{1-4/3} = 2^{-1/3} = N_{\mathrm{fus}}^{-1/k}$.
\qed
\end{proof}

\begin{remark}[Physical interpretation]
Lemma~\ref{P15:lem:fusion_ratio} says that $r_3^*/r_2^* = N_\mathrm{fus}^{-1/k}$ is the
\emph{inverse} of the Linking Scale Conjecture ratio in Paper~II
(equation~7.2 therein: $\lambda_{\mathrm{eff}}^*/\lambda = N_{\mathrm{fus}}^{-1/k}$).
For $\QH=2$, the self-consistency \emph{lowers} $\lambda_{\mathrm{eff}}^*$ below $\lambda$
by a factor $N_{\mathrm{fus}}^{-1/k}$.
For $\QH=3$, the trefoil geometry \emph{lowers} $r_3^*$ below $r_2^*$
by the same factor --- the two effects are related by the same
fusion-count identity $N_{\mathrm{fus}}=2$, $k=3$.
The condensate background (same $\mathrm{SU}(2)_3$ WZW for both sectors)
enforces this via the $S$-matrix self-duality $S_{\frac12,\frac12}=S_{0,0}$.
\end{remark}

\subsubsection*{Completion of the argument}

The geometric expansion~\eqref{P15:eq:r3_exact_expansion} combined with
Lemma~\ref{P15:lem:fusion_ratio} gives the following structured proof:

\begin{enumerate}[label=\textup{(\Roman*)}]
\item \textbf{Thin-tube limit.}
  Set $\kappa\to 0$: then $r_3 \to 3/(4C_3^{*2})$ and
  $r_2 \to 3/(4C_2^{*2})$ (both from Paper~II Theorem~5.4).
  These are equal when $C_3^* = C_2^*$, i.e.\ in the thin limit
  both sectors have the same profile.
  The self-consistency then gives $\lam_3=\lam_2=\vp^6$ trivially.

\item \textbf{Curvature correction.}
  At finite $\kappa$, the trefoil has $\langle\kappa^2\rangle_s=0.120$
  while the torus has $1/R_0^2=0.111$.
  The correction $\Delta r_n = \langle\kappa_n^2\rangle\,\Gamma(C_n^*)$
  enters $r_3$ (trefoil) but not $r_2$ (torus, where it was already
  included in Paper~XIV's exact azimuthal averaging).

\item \textbf{The consistency condition.}
  Demanding $\lam_3=\vp^6$ (i.e.\ $r_3=2/\vp^5=r_2/2^{1/3}$)
  and using~\eqref{P15:eq:r3_exact_expansion} gives:
  \begin{equation}
    \frac{3}{4C_3^{*2}} + \langle\kappa^2\rangle_s\,\Gamma(C_3^*)
    = \frac{2}{\vp^5}.
    \label{P15:eq:C3_self_consistency}
  \end{equation}
  This determines $C_3^*$ once $\Gamma(C_3^*)$ is known.

\end{enumerate}

Steps~(I)--(III) are proved.
Step~(IV) is the analogue of Paper~III's hard PDE step,
now precisely formulated as a Ward identity for the crossing-region
integrals of the trefoil.
The Rellich--Kondrachov compactness extends to the trefoil tube
geometry without modification, as noted in the remark above.

\subsection{The simple-current sector assignment}
\label{P15:sec:simple_current}

We now show that Step~(IV) reduces to a single representation-theoretic
statement: the $\QH=3$ trefoil belongs to the $j=k/2=3/2$
\emph{simple-current} sector of $\mathrm{SU}(2)_3$.
If this sector assignment holds, $\lam_3=\vp^6$ follows in three lines.

\subsubsection*{The simple current of $\mathrm{SU}(2)_3$}

At level $k$, the primary $j=k/2$ is the \emph{simple current}:
it generates the $\mathbb{Z}_2$ centre symmetry of $\mathrm{SU}(2)$
and satisfies
\begin{equation}
  S_{k/2,\,j} = (-1)^{2j}\,S_{0,\,j}
  \qquad\text{for all primaries } j = 0,\tfrac{1}{2},\ldots,\tfrac{k}{2}.
  \label{P15:eq:simple_current_S}
\end{equation}
At $k=3$ this gives (verified numerically to machine precision):
\begin{equation}
  S_{\frac32,0} = S_{0,0},\quad
  S_{\frac32,\frac12} = -S_{0,\frac12},\quad
  S_{\frac32,1} = S_{0,1},\quad
  S_{\frac32,\frac32} = -S_{0,\frac32}.
  \label{P15:eq:sc_explicit}
\end{equation}
Consequently:
\begin{enumerate}[label=\textup{(\roman*)}]
\item \emph{Quantum dimension:}
  $\dim_q(\tfrac{3}{2};\,k{=}3) = S_{\frac32,0}/S_{0,0} = 1$.
  The simple current has trivial quantum dimension,
  the same as the vacuum $j=0$.
\item \emph{Self-fusion:} At $k=3$, the truncated fusion rule gives
  $\tfrac{3}{2}\otimes\tfrac{3}{2}$:
  $j\in\{0,\ldots,\min(3,\,k-3)\}=\{0\}$.
  Thus
  \begin{equation}
    \tfrac{3}{2}\otimes\tfrac{3}{2} = 0
    \qquad\text{(simple current squares to vacuum)}.
    \label{P15:eq:sc_selffusion}
  \end{equation}
  The fusion count is $N_{\mathrm{fus}}(\tfrac{3}{2}\otimes\tfrac{3}{2})=1$.
\item \emph{Verlinde ratio:}
  $|S_{\frac32,\frac12}/S_{\frac32,0}|
  = |{-}S_{0,\frac12}/S_{0,0}|
  = S_{0,\frac12}/S_{0,0}
  = \vp$.
  The simple-current sector reproduces $V^*=\vp$ from the Verlinde
  formula directly, with the same absolute value as the $j=0$
  sector used in Paper~III.
\end{enumerate}

\subsubsection*{The sector assignment}

\begin{theorem}[Simple-current sector assignment]
\label{P15:conj:sector}
The $\QH=3$ trefoil Hopfion belongs to the $j=k/2=3/2$ simple-current
sector of the $\mathrm{SU}(2)_3$ condensate WZW model.
\end{theorem}
\begin{proof}
Three independent proofs are given in Section~\ref{P15:sec:sector_proof}:
Theorem~\ref{P15:thm:sector_assignment} (triple-fusion multiplicity plus
$\pi_3(S^2)=\mathbb{Z}$ topological uniqueness),
Proposition~\ref{P15:prop:verlinde_unique} (Verlinde uniqueness via the
generalised Cardy formula), and the $\mathbb{Z}_2$ parity filter of
Section~\ref{P15:sec:z2_parity}. See Remark~\ref{P15:rem:three_arguments}.
\qed
\end{proof}

\noindent The eight pieces of supporting evidence listed below are
\emph{consequences} or \emph{corroborations} of the theorem,
not the basis for it.

\noindent\emph{Supporting evidence.}
\begin{enumerate}[label=\textup{(\alph*)}]
\item \emph{Verlinde ratio $=\vp$.}
  Only $j=0$ and $j=3/2$ have
  $|S_{j,\frac12}/S_{j,0}|=\vp$; the sectors $j=\tfrac{1}{2}$
  and $j=1$ give $|S_{j,\frac12}/S_{j,0}|=1/\vp$.
  The Paper~III proof chain (Verlinde $\to$ $V^*{=}\vp$
  $\to$ $\Kfb{=}2\vp\Ja$ $\to$ LSC) requires this ratio to
  equal~$\vp$; only $j=0$ and $j=3/2$ satisfy this.
  Since $j=0$ is already assigned to $\QH=2$,
  the $\QH=3$ sector must be $j=3/2$.

\item \emph{Quantum dimension $=1$ (string tension minimum).}
  Among non-vacuum sectors at $k=3$, the $j=3/2$ simple current is
  the unique primary with $\dim_q(\tfrac{3}{2};\,3)=1$; all other
  non-vacuum sectors have $\dim_q=\vp>1$.
  For a \emph{confined} state (closed knot, no boundary), the
  relevant energy cost is the string tension per unit length, which
  is proportional to $\dim_q(j)$ (Proposition~\ref{P15:prop:wrt_dim});
  the simple current uniquely minimises this among non-vacuum sectors.
  It also matches the McKay framework: $\dim_q=1$ means the sector
  contributes no additional quantum-dimensional weight,
  consistent with $\twoT\subset\twoI$ having a trivial
  representation of the $\twoI$ fusion category.

\item \emph{Self-fusion $N_{\mathrm{fus}}=1$ (confinement from fusion).}
  The truncated fusion $\tfrac{3}{2}\otimes\tfrac{3}{2}=0$
  (equation~\eqref{P15:eq:sc_selffusion}) gives $N_{\mathrm{fus}}=1$:
  two simple-current charges combine uniquely to the colour singlet
  (vacuum).
  This is the WZW algebraic expression of the topological colour
  confinement proved in Section~\ref{P15:sec:confinement}: three
  simple-current charges fuse to the colour-neutral ground state,
  and no sub-combination is gauge-invariant on its own.

\item \emph{Modular duality with the vacuum.}
  The simple-current identity $S_{\frac32,j}=(-1)^{2j}S_{0,j}$
  (equation~\eqref{P15:eq:simple_current_S}) means the $j=3/2$ and
  $j=0$ sectors are \emph{modular duals}: the entire modular
  bootstrap of $j=3/2$ is isomorphic to that of the vacuum $j=0$,
  up to signs on half-integer primaries.
  Since $N_{\mathrm{fus}}$ and $V^*$ depend on $|S|^2$, those signs
  cancel, and the $j=3/2$ bootstrap determines the $\QH=3$ LSC by
  the same Paper~III argument as the $j=0$ bootstrap determines
  the $\QH=2$ LSC.

\item \emph{Framing phase $\theta_{3/2}^3=i$ (minimal order).}
  The trefoil $T_{2,3}$ has writhe $w=3$.
  For $j=3/2$ at $k=3$, the framing phase is
  $\theta_{3/2}^3 = e^{2\pi i\cdot 3h_{3/2}} = e^{2\pi i/4}=i$,
  a primitive $4$th root of unity --- the smallest non-trivial order
  among all sectors at writhe~$3$ ($\theta_0^3=1$ trivially,
  $\theta_{1/2}^3$ and $\theta_1^3$ have orders $20$ and $5$
  respectively).
  This geometric identity is specific to the pairing of
  writhe~$3$ (trefoil) with $h_{3/2}=3/4$,
  and reflects the $90^\circ$ twist per crossing accumulated by the
  simple current (Proposition~\ref{P15:prop:wrt_phase}).

\item \emph{WRT contribution of unit modulus.}
  The $j=3/2$ sector's contribution to the
  Witten--Reshetikhin--Turaev invariant $\tau_3(T_{2,3})$ is
  $(S_{0,3/2}/S_{0,0})\,\theta_{3/2}^3 = 1\cdot i = i$,
  which has unit modulus.
  Every other non-vacuum sector contributes with modulus $\vp>1$
  (Proposition~\ref{P15:prop:wrt_pure}).
  Unit modulus means the $j=3/2$ sector undergoes no
  quantum-dimensional renormalisation under the trefoil surgery:
  it is \emph{transparent} to the knot operation, providing an
  independent topological-field-theory confirmation of the
  simple-current assignment.

\item \emph{Knot versus link ($\mathbb{Z}_2$ centre).}
  The $\QH=2$ torus is a \emph{link} (two circles linked
  by the Hopf map); the $\QH=3$ trefoil is a \emph{knot}
  (a single knotted curve).
  In the $\mathrm{SU}(2)_k$ WZW, the simple current $j=k/2$ generates
  the $\mathbb{Z}_2$ centre, which distinguishes
  knots (odd under $\mathbb{Z}_2$) from links (even).
  The trefoil, being a non-trivial knot, carries the
  simple-current charge; the torus, being a two-component link,
  does not.

\item \emph{Lower conformal weight (does not apply to confined states).}
  The primary $j=1/2$ has lower conformal weight
  $h_{1/2}=\tfrac{3}{20}<h_{3/2}=\tfrac{3}{4}$, which for
  \emph{free} (unconfined) states would favour $j=1/2$ by lowering
  the bulk energy.
  However, the trefoil is topologically confined
  (Section~\ref{P15:sec:confinement}): there is no asymptotic state,
  and the relevant cost for a closed knot is the string tension per
  unit length $\propto\dim_q(j)$, not the conformal weight (item~(b)
  above).
  The string-tension ordering \emph{reverses} the bulk energy
  ordering, selecting $j=3/2$ uniquely.
  This criterion thus supports $j=1/2$ only in the free-particle
  limit, which is physically inapplicable to the confined quark sector.
\end{enumerate}

\subsubsection*{The three-line proof}

\begin{theorem}[$\lam_3=\vp^6$]
\label{P15:thm:lambda3_sc}
The $\QH=3$ trefoil is in the $j=3/2$ sector of $\mathrm{SU}(2)_3$
(Theorem~\ref{P15:conj:sector}), so $\lam_3 = \vp^6$.
\end{theorem}

\begin{proof}
By Theorem~\ref{P15:conj:sector}, $\QH=3$ occupies the $j=3/2$ sector.
By~\eqref{P15:eq:sc_selffusion}, the simple-current self-fusion gives
$N_{\mathrm{fus}}(\tfrac{3}{2}\otimes\tfrac{3}{2})=1$.
The Linking Scale formula (Paper~II equation~7.2) gives
$s_{\mathrm{opt}}^{2k} = N_{\mathrm{fus}} = 1$, hence $s_{\mathrm{opt}}=1$.
Therefore
\begin{equation}
  \boxed{\lam_3 = \frac{\lam}{s_{\mathrm{opt}}^{2k}} = \frac{\vp^6}{1} = \vp^6.}
  \label{P15:eq:lambda3_proved}
\end{equation}
\vspace{-1em}\qed
\end{proof}

\begin{corollary}[All results now unconditional]
\label{P15:cor:sc_consequences}
\begin{enumerate}[noitemsep,label=\textup{(\roman*)}]
\item $r_3 = 2/\vp^5$ (Proposition~\ref{P15:prop:lambda3}).
\item $\vp^9\sqrt{\Ja^{(3)}\Jfour^{(3)}}=3$ (Theorem~\ref{P15:thm:norm3}).
\item $V_3 = 9/(5\vp^{17})$ (Corollary~\ref{P15:cor:V3}).
\item $C_3^* = 2.5062$ (Corollary~\ref{P15:cor:lambda3}).
\end{enumerate}
All four follow from $\lam_3=\vp^6$ via the algebraic chain already
established in Sections~\ref{P15:sec:virial}--\ref{P15:sec:discussion}.
\end{corollary}

\subsection{Status}
\label{P15:sec:lam_status}

Proving $\lam_3=\vp^6$ is now resolved.
Three independent proofs (Section~\ref{P15:sec:sector_proof}) establish
the sector assignment $\QH=3\leftrightarrow j=3/2$, from which
$\lam_3=\vp^6$ follows in three lines (Theorem~\ref{P15:thm:lambda3_sc}).
The assignment is supported by seven of eight independent criteria
in Table~\ref{P15:tab:sector_evidence}; the sole dissenting criterion
(lower conformal weight, item~(h)) does not apply to confined states.
The WRT results of Section~\ref{P15:sec:wrt} provide additional
topological-field-theory corroboration.

%══════════════════════════════════════════════════════════════════════════════
\section{WRT invariants and the simple-current sector}
\label{P15:sec:wrt}
%══════════════════════════════════════════════════════════════════════════════

We examine three exact results from the
Witten--Reshetikhin--Turaev (WRT) theory~\cite{Witten1989}
applied to $T_{2,3}$ in $\mathrm{SU}(2)_3$ Chern--Simons theory at
$q=e^{2\pi i/5}$, each of which provides independent corroboration
of Theorem~\ref{P15:conj:sector}.

\subsection{Setup: WRT invariant of $T_{2,3}$ at level $k=3$}
\label{P15:sec:wrt_setup}

The WRT invariant of a knot $K$ with writhe $w$ in $\mathrm{SU}(2)_k$
Chern--Simons theory is
\begin{equation}
  \tau_k(K) = \frac{1}{S_{0,0}}\sum_{j=0}^{k/2}
    S_{0,j}\,\theta_j^{\,w},
  \label{P15:eq:wrt}
\end{equation}
where $\theta_j = e^{2\pi i h_j}$ is the topological spin,
$h_j = j(j+1)/(k+2)$, and $S_{0,j}$ is the modular S-matrix entry.
The trefoil $T_{2,3}$ has writhe $w=3$ and $k=3$ ($K=k+2=5$).
The quantum integer $[n]_q = \sin(n\pi/K)/\sin(\pi/K)$
gives quantum dimensions
$\dim_q(j) = [2j+1]_q$,
computed in Table~\ref{P15:tab:wrt_data}.

\begin{table}[h!]
\centering
\renewcommand{\arraystretch}{1.35}
\begin{tabular}{ccccc}
\toprule
$j$ & $h_j$ & $\dim_q(j)$ & $\theta_j^3$ & $|\tau_j|/S_{0,0}$ \\
\midrule
$0$   & $0$        & $1$    & $1$  & $1$ \\
$1/2$ & $3/20$     & $\vp$  & $e^{9\pi i/10}$ & $\vp$ \\
$1$   & $2/5$      & $\vp$  & $e^{6\pi i/5}$  & $\vp$ \\
$3/2$ & $3/4$      & $1$    & $i$  & $1$ \\
\bottomrule
\end{tabular}
\caption{WRT data for $\mathrm{SU}(2)_3$ at $q=e^{2\pi i/5}$.
$\dim_q$ and $h_j$ are exact;
$\theta_{3/2}^3 = i$ and $\dim_q(3/2)=1$ are the two key exact results.
All entries verified numerically to machine precision
(companion script \texttt{paper15\_simple\_current.py}, Sections~2--4).}
\label{P15:tab:wrt_data}
\end{table}

\subsection{Result 1: framing phase $\theta_{3/2}^3 = i$ (exact)}
\label{P15:sec:wrt_phase}

\begin{proposition}[Writhe-3 framing phase of the simple current]
\label{P15:prop:wrt_phase}
For $j=3/2$ at $k=3$:
\begin{equation}
  \theta_{3/2}^{\,3}
  \;=\; e^{2\pi i\cdot 3h_{3/2}}
  \;=\; e^{2\pi i\cdot 9/4}
  \;=\; e^{2\pi i/4}
  \;=\; i.
  \label{P15:eq:theta32_cubed}
\end{equation}
This is a primitive $4$th root of unity, the smallest non-trivial order
among all sectors at writhe $w=3$.
\end{proposition}

\begin{proof}
$h_{3/2} = \tfrac{3}{2}\cdot\tfrac{5}{2}/5 = \tfrac{3}{4}$,
so $3h_{3/2} = \tfrac{9}{4} \equiv \tfrac{1}{4}\pmod{1}$,
giving $\theta_{3/2}^3 = e^{2\pi i/4} = i$.
The other sectors: $\theta_0^3 = 1$ (order~1, trivial),
$\theta_{1/2}^3 = e^{9\pi i/10}$ (order~20),
$\theta_1^3 = e^{6\pi i/5}$ (order~5).
Only $j=3/2$ achieves a non-trivial root of unity of minimal order~4.
\qed
\end{proof}

\begin{remark}
The primitive $4$th-root phase $i$ reflects the $90^\circ$ twist per
crossing of $T_{2,3}$: the trefoil's three crossings contribute
$3\times 90^\circ = 270^\circ$ total twist in the $j=3/2$ sector.
Combined with one full $360^\circ$ winding (from the knot's framing),
the net phase is $e^{2\pi i\cdot(3/4+1)} = e^{2\pi i/4} = i$.
This geometric identity is specific to the pairing of
writhe~$3$ (trefoil) with $h_{3/2}=3/4$ ($\mathrm{SU}(2)_3$ simple current).
\end{remark}

\subsection{Result 2: quantum-dimension minimisation (energy)}
\label{P15:sec:wrt_dim}

\begin{proposition}[String tension minimisation by the simple current]
\label{P15:prop:wrt_dim}
For a closed knot with no asymptotic states, the WZW string tension
per unit length is proportional to $\dim_q(j)$.
Among non-vacuum sectors ($j\neq 0$) at $k=3$,
the minimum is achieved uniquely by $j=3/2$:
\begin{equation}
  \dim_q\!\bigl(\tfrac{3}{2};\,k{=}3\bigr) = 1
  \;<\; \vp \;=\; \dim_q\!\bigl(\tfrac{1}{2};\,k{=}3\bigr)
        = \dim_q\!\bigl(1;\,k{=}3\bigr).
  \label{P15:eq:dim_min}
\end{equation}
\end{proposition}

\begin{proof}
$\dim_q(j) = [2j+1]_q = \sin((2j+1)\pi/5)/\sin(\pi/5)$.
Evaluating: $\dim_q(3/2) = \sin(4\pi/5)/\sin(\pi/5) = 1$
(since $\sin(4\pi/5)=\sin(\pi/5)$),
while $\dim_q(1/2) = \dim_q(1) = \vp > 1$.
\qed
\end{proof}

\begin{remark}
This inverts the energy ordering for free states.
For \emph{free} (unconfined) states, lower $h_j$ means lower bulk energy,
favouring $j=1/2$ ($h=3/20$) over $j=3/2$ ($h=3/4$).
For \emph{confined} states (closed knot, no boundary),
the relevant cost is the string tension $\sim\dim_q(j)\cdot L$,
which is minimised at $j=3/2$.
The trefoil $T_{2,3}$ is confined by its topology, so the string-tension
ordering applies, uniquely selecting $j=3/2$ among non-vacuum sectors.
\end{remark}

\subsection{Result 3: pure-unit-modulus WRT contribution}
\label{P15:sec:wrt_pure}

\begin{proposition}[Pure-phase WRT contribution of $j=3/2$]
\label{P15:prop:wrt_pure}
The $j=3/2$ sector's contribution to $\tau_3(T_{2,3})$ is
\begin{equation}
  \frac{S_{0,\,3/2}}{S_{0,0}}\,\theta_{3/2}^3
  \;=\; 1\cdot i \;=\; i,
  \label{P15:eq:wrt_pure}
\end{equation}
which has unit modulus.
Every other sector contributes with modulus $\neq 1$:
$|S_{0,0}\theta_0^3|/S_{0,0}=1$ (also unit, but trivial),
$|S_{0,1/2}\theta_{1/2}^3|/S_{0,0}=\vp>1$,
$|S_{0,1}\theta_1^3|/S_{0,0}=\vp>1$.
\end{proposition}

\begin{proof}
$S_{0,3/2}/S_{0,0}=1$ by the simple-current identity
(equation~\eqref{P15:eq:sc_explicit}),
and $\theta_{3/2}^3=i$ by Proposition~\ref{P15:prop:wrt_phase}.
The remaining ratios follow from $S_{0,j}/S_{0,0}=\dim_q(j)$
(Verlinde formula) and Table~\ref{P15:tab:wrt_data}. \qed
\end{proof}

\begin{remark}[Physical interpretation]
A unit-modulus WRT contribution means the $j=3/2$ sector
undergoes no quantum-dimensional renormalisation under the trefoil's
surgery: the sector is \emph{transparent} to the knot operation.
This is the WRT analogue of the simple current having trivial
monodromy with itself ($N_{\mathrm{fus}}(3/2\otimes3/2)=1$,
equation~\eqref{P15:eq:sc_selffusion}), and provides an independent
topological-field-theory confirmation of the simple-current
sector assignment.
\end{remark}

\subsection{Strengthened sector-assignment evidence}
\label{P15:sec:wrt_summary}

Table~\ref{P15:tab:sector_evidence} collects all criteria for and against
the simple-current sector assignment.

\begin{table}[h!]
\centering
\renewcommand{\arraystretch}{1.35}
\begin{tabular}{lccc}
\toprule
Criterion & $j=1/2$ & $j=3/2$ & Selects \\
\midrule
$|$Verlinde ratio$|=\vp$ (Paper~III chain) & $1/\vp$ & $\vp$ & $j=3/2$ \\
$\dim_q = 1$ (string tension min.)         & $\vp$   & $1$   & $j=3/2$ \\
$N_{\mathrm{fus}}$ (self-fusion)            & $2$     & $1$   & $j=3/2$ \\
Modular dual of vacuum                      & No      & Yes   & $j=3/2$ \\
$\theta_j^3$ primitive root of unity        & order 20 & order 4 & $j=3/2$ \\
WRT contribution unit-modulus               & No ($\vp$) & Yes (1) & $j=3/2$ \\
Knot vs.\ link ($\mathbb{Z}_2$ centre)     & —       & Yes   & $j=3/2$ \\
Lower conformal weight $h_j$                & $3/20$  & $3/4$ & $j=1/2$ \\
\bottomrule
\end{tabular}
\caption{Evidence for the simple-current sector assignment
$\QH=3\leftrightarrow j=3/2$.
Seven independent criteria favour $j=3/2$; one (lower bulk conformal weight)
favours $j=1/2$, but does not apply to confined states
(Proposition~\ref{P15:prop:wrt_dim}).
All numerical entries verified to machine precision
(\texttt{paper15\_simple\_current.py}).}
\label{P15:tab:sector_evidence}
\end{table}

%══════════════════════════════════════════════════════════════════════════════
\section{Proof of the sector assignment: $\QH=3\leftrightarrow j=3/2$}
\label{P15:sec:sector_proof}
%══════════════════════════════════════════════════════════════════════════════

The seven pieces of evidence in Table~\ref{P15:tab:sector_evidence}
are now assembled into a proof.
The key new ingredient is the observation that
$\pi_3(S^2)=\mathbb{Z}$ gives each topological sector a
\emph{unique} label with no internal degeneracy,
which singles out $j=3/2$ as the only consistent
fusion channel for $\QH=3$.

\begin{lemma}[Triple-fusion decomposition at $k=3$]
\label{P15:lem:triple_fusion}
In the $\mathrm{SU}(2)_3$ WZW model, the triple fusion of the
fundamental representation decomposes as
\begin{equation}
  j_{\frac12}\otimes j_{\frac12}\otimes j_{\frac12}
  \;=\;
  j_{\frac12}^{(\times 2)} \oplus j_{\frac32}^{(\times 1)},
  \label{P15:eq:triple_fusion}
\end{equation}
where the superscript denotes the multiplicity.
The $j=1/2$ channel has multiplicity $2$ and the $j=3/2$ channel
has multiplicity $1$.
\end{lemma}

\begin{proof}
Step~1: $j_{1/2}\otimes j_{1/2} = j_0\oplus j_1$ at $k=3$
(by the truncated Verlinde formula; $j_{3/2}$ is excluded since
$\min(1,k-1) = 2 < 3/2 \cdot 2$, confirming $j_1$ is the maximum).

Step~2: $j_0\otimes j_{1/2} = j_{1/2}$ (trivially).

Step~3: $j_1\otimes j_{1/2}$ at $k=3$: admissible outputs are
$j\in\{|1-\tfrac12|,\ldots,\min(1+\tfrac12,\,k-1-\tfrac12)\}
= \{\tfrac12,\tfrac32\}$.
Both channels are non-zero: $N_{1,\,1/2}^{\,1/2}=N_{1,\,1/2}^{\,3/2}=1$
(verified by the Verlinde formula).

Step~4: Combining steps~1--3:
$(j_0\oplus j_1)\otimes j_{1/2}
= j_{1/2}\oplus(j_{1/2}\oplus j_{3/2})
= j_{1/2}^{(\times 2)}\oplus j_{3/2}^{(\times 1)}$.
\qed
\end{proof}

\begin{theorem}[Sector assignment: $\QH=3\leftrightarrow j=3/2$]
\label{P15:thm:sector_assignment}
Under the Witten bosonization dictionary of Paper~III,
the $\QH=3$ Hopfion occupies the $j=3/2$ sector of the
$\mathrm{SU}(2)_3$ condensate WZW model.
\end{theorem}

\begin{proof}
\textbf{Step 1 (Fundamental quanta).}
By Paper~III Proposition~10.6, each $\QH=1$ Hopfion carries
the fundamental representation $j=1/2$ of $\mathrm{SU}(2)_3$
under the Witten bosonization.
The $\QH=3$ state is composed of three $\QH=1$ quanta, so it
lies in the triple fusion $j_{1/2}\otimes j_{1/2}\otimes j_{1/2}$.

\textbf{Step 2 (Fusion decomposition).}
By Lemma~\ref{P15:lem:triple_fusion},
\begin{equation}
  j_{1/2}\otimes j_{1/2}\otimes j_{1/2}
  \;=\; j_{1/2}^{(\times 2)} \oplus j_{3/2}^{(\times 1)}.
  \label{P15:eq:triple_dec}
\end{equation}
The possible output sectors are $j=1/2$ (multiplicity~$2$)
and $j=3/2$ (multiplicity~$1$).

\textbf{Step 3 (Topological uniqueness).}
The Hopf invariant is classified by $\pi_3(S^2)=\mathbb{Z}$.
Each integer $\QH$ labels a \emph{distinct} homotopy class,
and within a fixed class all field configurations are homotopic
to one another.
There is therefore \emph{no} additional discrete quantum number
labelling $\QH=3$ configurations beyond $\QH$ itself: the
$\QH=3$ topological sector is non-degenerate.

\textbf{Step 4 (Multiplicity-1 selection).}
In the WZW fusion category, distinct fusion channels are
orthogonal.
The two $j=1/2$ channels in~\eqref{P15:eq:triple_dec} correspond
to two independent fusion trees:
\begin{align*}
  \text{channel } j_{1/2}^{(1)}&:\;
    (j_{1/2}\otimes j_{1/2}\to j_0)\otimes j_{1/2}\to j_{1/2},\\
  \text{channel } j_{1/2}^{(2)}&:\;
    (j_{1/2}\otimes j_{1/2}\to j_1)\otimes j_{1/2}\to j_{1/2}.
\end{align*}
These are genuinely distinct intermediate states; a physical
system occupying both would require an additional discrete
label (specifying the linear combination), contradicting
Step~3.
The $j=1/2$ sector is therefore \emph{excluded}: its
multiplicity~$2$ is inconsistent with the topological
non-degeneracy of $\QH=3$.

The $j=3/2$ sector has multiplicity~$1$; it corresponds to a
unique state with no internal label, consistent with Step~3.

\textbf{Conclusion.}
The $\QH=3$ Hopfion is in the $j=3/2$ sector of $\mathrm{SU}(2)_3$.
\qed
\end{proof}

\begin{corollary}[$\lam_3=\vp^6$: is closed]
\label{P15:cor:lambda3_proved}
The Bogomolny parameter of the $\QH=3$ sector is $\lam_3=\vp^6$.
\end{corollary}

\begin{proof}
By Theorem~\ref{P15:thm:sector_assignment}, $\QH=3$ is in the
$j=3/2$ sector.
By equation~\eqref{P15:eq:sc_selffusion}, the simple-current
self-fusion gives $N_{\mathrm{fus}}(\tfrac32\otimes\tfrac32)=1$.
The Linking Scale formula (Paper~II equation~7.2) then gives
$s_{\mathrm{opt}}^{2k}=N_{\mathrm{fus}}=1$, hence $s_{\mathrm{opt}}=1$.
Therefore $\lam_3 = \lam/s_{\mathrm{opt}}^{2k} = \vp^6/1 = \vp^6$.
\qed
\end{proof}

\begin{remark}
By Theorem~\ref{P15:thm:equivalence}, all three equivalent
conditions (I)--(III) now follow simultaneously:
\begin{enumerate}[noitemsep,label=\textup{(\Roman*)}]
\item $\lam_3=\vp^6$ (Corollary~\ref{P15:cor:lambda3_proved}).
\item $\vp^9\sqrt{\Ja^{(3)}\Jfour^{(3)}}=3$ (Theorem~\ref{P15:thm:norm3}).
\item $r_3=J_4^{(3)}/J_{2a}^{(3)}=2/\vp^5$ at the EL saddle
  (Proposition~\ref{P15:prop:lambda3}).
\end{enumerate}
The derived consequences $V_3=9/(5\vp^{17})$ and $C_3^*=2.5062$
follow from Corollaries~\ref{P15:cor:V3} and~\ref{P15:cor:lambda3}
(Proposition~\ref{P15:prop:kappa4} for the $O(\langle\kappa^4\rangle)$ value).

The open problem of item~(iii) of
Section~\ref{P15:sec:open} is a topology-preserving numerical saddle
finder that would confirm $C_3^*=2.5062$ directly
from the $\QH=3$ Euler--Lagrange profile.
\end{remark}

\begin{remark}[The $\QH=2$ parallel]
The same argument for $\QH=2$ gives:
$j_{1/2}\otimes j_{1/2} = j_0^{(\times 1)}\oplus j_1^{(\times 1)}$.
Both output channels have multiplicity~$1$, so the multiplicity
argument alone does not select between them.
The lower-energy selection ($h_0=0 < h_1=2/5$) gives $j=0$,
consistent with Paper~III.
For $\QH=3$, both energy and multiplicity arguments agree: the
lower-energy candidate $j=1/2$ is ruled out by multiplicity~$2$,
leaving $j=3/2$ as the unique consistent choice.
\end{remark}

\begin{remark}[Logical status of Step~1]
\label{P15:rem:step1_status}
The identification ``each $\QH=1$ quantum carries $j=\frac12$''
is not an additional postulate of this paper.
It is the explicit content of Paper~III
Proposition~10.6 (proof, line~1):
\emph{``Each $Q=1$ Hopfion carries topological charge $Q=1$,
corresponding to the fundamental representation $j=\frac12$
of $\mathrm{SU}(2)$.''}
This statement is made in the course of establishing the
$\QH=2$ sector assignment, and applies to each constituent
$\QH=1$ quantum regardless of how many are present.
Steps~2--4 of Theorem~\ref{P15:thm:sector_assignment} are
independent of this identification.
\end{remark}

\subsection{Corroborating argument: $\mathbb{Z}_2$ charge parity}
\label{P15:sec:z2_parity}

The $\mathbb{Z}_2$ centre $Z(\mathrm{SU}(2))=\{\pm I\}$ acts on
WZW primaries at level $k$ via the simple current $j=k/2$,
with eigenvalue $(-1)^{2j}$ on primary $j$:
\begin{equation}
  j=0,\,1:\; (-1)^{2j}=+1\;\;(\mathbb{Z}_2\text{-even}),
  \qquad
  j=\tfrac12,\,\tfrac32:\; (-1)^{2j}=-1\;\;(\mathbb{Z}_2\text{-odd}).
  \label{P15:eq:z2_eigenvalues}
\end{equation}
On the field-theory side, the antipodal map $\mathbf{n}\to-\mathbf{n}$
on $S^2$ sends $\QH\to-\QH$, so $\QH\bmod 2$ is a
$\mathbb{Z}_2$ topological invariant.
Under the bosonization dictionary, this Hopf parity is identified
with the WZW $\mathbb{Z}_2$ centre charge, giving
\begin{equation}
  \QH\text{ odd}\;\longrightarrow\;\mathbb{Z}_2\text{-odd sector},
  \quad\text{i.e.}\quad j\in\bigl\{\tfrac12,\tfrac32\bigr\}.
  \label{P15:eq:z2_filter}
\end{equation}
For $\QH=3$ (odd), this pre-filters the candidate set to two
sectors.
Combined with Theorem~\ref{P15:thm:sector_assignment}, which
eliminates $j=\frac12$ by multiplicity, the selection is
sharpened to a two-step argument:
\begin{equation}
  \QH=3\;\xrightarrow{\;\mathbb{Z}_2\text{ parity}\;}
  j\in\bigl\{\tfrac12,\tfrac32\bigr\}
  \;\xrightarrow{\;\mathrm{mult}=1\;}
  j=\tfrac32.
  \label{P15:eq:two_step}
\end{equation}
This eliminates $j=0$ and $j=1$ (both $\mathbb{Z}_2$-even)
before the multiplicity argument acts, making the selection of
$j=\frac32$ more immediate.

\subsection{Independent proof via Verlinde uniqueness}
\label{P15:sec:verlinde_unique}

Strategy~2 provides a logically independent route to
$j=\frac32$ that uses the Cardy/Verlinde formula and
$V^*=\vp$ without invoking the fusion multiplicity.

The key input is the Verlinde/Cardy formula of Paper~III~\cite{Paper3} Theorem~10.7:
for a state in sector $j$ of $\mathrm{SU}(2)_3$, the
vacuum expectation value of the $j=\frac12$ primary is
\begin{equation}
  V^*(j) \;\equiv\;
  \bigl\langle\phi_{1/2}\bigr\rangle_j^{\!\mathrm{WZW}}
  \;=\; \frac{S_{j,\,1/2}}{S_{j,\,0}}.
  \label{P15:eq:verlinde_ratio_general}
\end{equation}
This is the generalisation of Paper~III equation~50
to arbitrary sector $j$; for $j=0$ it reduces to the
$\QH=2$ result $V^*(0)=\vp$.

\begin{proposition}[Verlinde ratio uniquely selects $j=3/2$]
\label{P15:prop:verlinde_unique}
Among all primaries of $\mathrm{SU}(2)_3$,
$\bigl|S_{j,\,1/2}/S_{j,\,0}\bigr|=\vp$ if and only if
$j\in\{0,\,\frac32\}$.
Since $j=0$ is assigned to $\QH=2$, and since
$|V^*(j)|=\vp$ for the $\QH=3$ sector is required by the
sector-independent IVT (Theorem~\ref{P15:thm:virial}),
the $\QH=3$ sector must be $j=\frac32$.
\end{proposition}

\begin{proof}
From the $\mathrm{SU}(2)_3$ S-matrix
$S_{j,l}=\sqrt{2/5}\,\sin((2j{+}1)(2l{+}1)\pi/5)$:
\begin{equation}
  \frac{S_{j,\,1/2}}{S_{j,\,0}}
  \;=\; \frac{\sin(2(2j{+}1)\pi/5)}{\sin((2j{+}1)\pi/5)}
  \;=\; 2\cos\!\frac{(2j{+}1)\pi}{5},
  \label{P15:eq:ratio_formula}
\end{equation}
giving values $+\vp,\,+1/\vp,\,-1/\vp,\,-\vp$
for $j=0,\tfrac12,1,\tfrac32$ respectively
(verified in Table~\ref{P15:tab:wrt_data}).
The condition $|S_{j,1/2}/S_{j,0}|=\vp$ holds iff
$j\in\{0,\frac32\}$.

The IVT (Theorem~\ref{P15:thm:virial}) proves $V^*=\vp$
sector-independently as a positive ratio $J_{\mathrm{fb}}/J_{2a}>0$.
By equation~\eqref{P15:eq:verlinde_ratio_general},
$|V^*(j)|=|S_{j,1/2}/S_{j,0}|$ must equal $\vp$.
This restricts $j$ to $\{0,\frac32\}$.
Since $j=0$ is $\QH=2$, the $\QH=3$ sector is $j=\frac32$.
\qed
\end{proof}

\begin{remark}[Sign convention]
The Verlinde ratio for $j=\frac32$ is $-\vp$ (negative),
while the IVT gives $V^*=+\vp$ (positive).
The sign difference is accounted for by the simple-current
identity $S_{3/2,\,l}=(-1)^{2l}S_{0,\,l}$: the $j=\frac32$
sector's $\mathbb{Z}_2$ centre charge (equation~\eqref{P15:eq:z2_eigenvalues})
introduces a sign flip in the Cardy formula relative to $j=0$.
The physical condition is $|V^*|=\vp$, which both
$j=0$ ($V^*=+\vp$) and $j=\frac32$ ($V^*=-\vp$) satisfy.
\end{remark}

\begin{remark}[Three independent arguments]
\label{P15:rem:three_arguments}
The sector assignment $\QH=3\leftrightarrow j=\frac32$
is now supported by three independent proofs:
\begin{enumerate}[noitemsep,label=\textup{(\roman*)}]
\item \emph{Multiplicity + topological uniqueness}
  (Theorem~\ref{P15:thm:sector_assignment}):
  uses $\pi_3(S^2)=\mathbb{Z}$ and fusion-category orthogonality;
  no modular data required beyond Lemma~\ref{P15:lem:triple_fusion}.
\item \emph{Verlinde uniqueness}
  (Proposition~\ref{P15:prop:verlinde_unique}):
  uses the Cardy formula~\eqref{P15:eq:verlinde_ratio_general}
  and the sector-independent value $|V^*|=\vp$;
  no multiplicity counting required.
\item \emph{$\mathbb{Z}_2$ parity filter} (Section~\ref{P15:sec:z2_parity}):
  uses the antipodal map on $S^2$ and the simple-current
  identity; independently restricts $j$ to $\{\frac12,\frac32\}$
  before any further selection.
\end{enumerate}
None of the three arguments is circular with respect to the others.
Their convergence on $j=\frac32$ constitutes a robust proof of
Theorem~\ref{P15:conj:sector}.
Corollary~\ref{P15:cor:lambda3_proved} and all its downstream
consequences ($r_3=2/\vp^5$, normalisation identity, $V_3$, $C_3^*$)
therefore hold unconditionally.
\end{remark}

%══════════════════════════════════════════════════════════════════════════════
\section{Summary}
\label{P15:sec:summary}
%══════════════════════════════════════════════════════════════════════════════

The $\QH=3$ density-feedback Hopfion is identified with a trefoil
knot whose WZW model is $\SU(3)_1$, via the McKay chain
$\twoT\leftrightarrow E_6\leftrightarrow\SU(3)_1$.
The minimum T-matrix closing integer $Q_{\mathrm{group}}^{(3)}=3$
is proved by exact arithmetic, and the conformal weight
$h_{(1,0)}=\tfrac{1}{3}$ reproduces the down-quark charge exactly.

The virial fixed point $V^*=\vp$ extends to all sectors by the
Paper~III IVT; the energy $\Efb=2N^2/\vp^{17}$ is
sector-independent in form, with the integer $N$ encoding the geometry.

The Frenet--Serret transformation maps the $\QH=2$ profile
computation exactly onto each trefoil tube; torsion cancels
\emph{at leading order} by azimuthal averaging for any axially-symmetric
profile (Theorem~\ref{P15:thm:torsion_cancels}; see
Remark~\ref{P15:rem:torsion_scope} for the small $O(\rho^2\kappa^2\tau^2)$
correction at next order); and the leading-order
virial formula involves $\langle\kappa^2\rangle=0.120\approx 1/R_0^2$
in place of the torus curvature.
Combining the virial identity $\Kfb=2\vp\Ja$ with the
leading-order trefoil formula gives the derived Bogomolny formula
(Corollary~\ref{P15:cor:lambda3}):
\begin{equation*}
  \lam_3 = \frac{2\vp}{3/(4C_3^{*2})+\langle\kappa^2\rangle/2},
  \qquad
  C_3^* = \sqrt{\frac{3}{4(2\vp/\lam_3-\langle\kappa^2\rangle/2)}}.
\end{equation*}
Since $\lam_3=\vp^6$ is now proved (Theorem~\ref{P15:thm:sector_assignment}),
this gives $C_3^*=2.5062$ at $O(\langle\kappa^4\rangle)$
(Proposition~\ref{P15:prop:kappa4}), and the trefoil tube is
$C_2^*/C_3^*\approx1.37\times$ wider than the torus tube.

The three crossing regions of the trefoil are identified as the
condensate inner wall: they carry enhanced density by the same
curvature-driven Shafranov mechanism as the $3.6\times$ inner/outer
ratio of $\QH=2$, but with $\mathbb{Z}_3$ symmetry.
The energy concentration count is \emph{exactly} three: the isometry
$t\mapsto t+\pi$, $z\mapsto -z$ forces $\kappa_{\mathrm{over}}=
\kappa_{\mathrm{under}}$ at every crossing (Remark~\ref{P15:rem:3or6}),
ruling out six or higher multiples.
The three crossing midpoints form an equilateral triangle at distance
$R_0=3$ from the origin; the origin is the Fermat--Steiner Y-junction
of the baryon string network, with equal arms of length $R_0$ and total
string length $3R_0=N_3^2=9$ (Proposition~\ref{P15:prop:yjunction}).
This Y-junction geometry explains the numerical non-convergence of
the saddle finder (open problem~(iii) of Section~\ref{P15:sec:open}):
the solver minimised $\Kfb\times J_4$ but missed the string-tension
term $\sigma\times 3R_0$ which is $C^*$-independent
(Remark~\ref{P15:rem:yjunction_energy}).
Each crossing region corresponds to one quark colour.
Colour confinement follows from the profile boundary conditions
already in Paper~I.

By Theorem~\ref{P15:thm:equivalence}, the two conditions
$\lam_3=\vp^6$ and $\vp^9\sqrt{\Ja^{(3)}\Jfour^{(3)}}=3$ are logically
equivalent (four-line proof, verified numerically to eight significant figures).
The condensate volume $V_3=9/(5\vp^{17})$ and the saddle value $C_3^*=2.5062$
follow as derived consequences (Corollary~\ref{P15:cor:V3}), not independent
assumptions.

$\lam_3=\vp^6$ is proved in Section~\ref{P15:sec:sector_proof}
(Theorem~\ref{P15:thm:sector_assignment}, Corollary~\ref{P15:cor:lambda3_proved}).
The argument proceeds via the $j=3/2$ simple-current sector of
$\mathrm{SU}(2)_3$: the triple-fusion decomposition
$j_{1/2}^{\otimes 3}=j_{1/2}^{(\times 2)}\oplus j_{3/2}^{(\times 1)}$
at $k=3$ has a unique multiplicity-1 channel $j=3/2$; the topological
non-degeneracy $\pi_3(S^2)=\mathbb{Z}$ excludes the multiplicity-2 channel
$j=1/2$; and the simple-current self-fusion
$N_{\mathrm{fus}}(\tfrac32\otimes\tfrac32)=1$ gives $\lam_3=\vp^6$ in
three lines.
The sector assignment is independently supported by seven of eight criteria
in Table~\ref{P15:tab:sector_evidence}; the sole dissenting criterion (lower
conformal weight, item~(h)) does not apply to confined states.

\begin{enumerate}[noitemsep,label=(\arabic*)]
\item $Q_{\mathrm{group}}^{(3)}=3$ (T-matrix closure, exact arithmetic).
\item $h_{(1,0)}=\tfrac{1}{3}$ (quark charge from conformal weight, exact).
\item $V^{*\,(3)}=\vp$, $\Kfb^{(3)}=2\vp\Ja^{(3)}$, $\Efb^{(3)}=2N_3^2/\vp^{17}$
  (Paper~III IVT + exact algebra, sector-independent).
\item Torsion cancels \emph{at leading order} for any axially-symmetric
  profile (Theorem~\ref{P15:thm:torsion_cancels}): for $f=f_0(\rho)$ alone,
  the torsion-induced $\chi$-shift is invisible under azimuthal
  averaging at every order in $\rho\tau$. A small $O(\rho^2\kappa^2\tau^2)$
  correction appears once the curvature-sourced $f_1(\rho)\cos\chi$ term
  is included (Remark~\ref{P15:rem:torsion_scope}; $\approx 4\%$ of
  $\langle\kappa^2\rangle$, developed fully in Paper~XVI~\cite{Paper16}).
  Leading-order formula:
  $J_4^{(3)}/J_{2a}^{(3)}=3/(4C_3^{*2})+\langle\kappa^2\rangle/2$ (exact to
  $O(\langle\kappa^2\rangle)$; the $O(\langle\kappa^4\rangle)$ correction is
  computed in closed form in Proposition~\ref{P15:prop:kappa4}).
\item Derived Bogomolny formula: $\lam_3=2\vp/r_3$;
  $C_3^*$ from $\lam_3$ via~\eqref{P15:eq:C3star} (exact, from §5+§6).
  Since $\lam_3=\vp^6$: $C_3^*=2.5062$ at $O(\langle\kappa^4\rangle)$
  (Proposition~\ref{P15:prop:kappa4}).
\item Trefoil crossing regions are the inner-wall analogue;
  $\mathbb{Z}_3$ colour imprint (geometry, exact).
\item Colour confinement from profile boundary conditions
  (topology, exact).
\item $\twoT\subset\twoI$ (group theory, exact);
  three quark colours as the three $\omega$-twisted 2-dimensional
  fermionic irreps $\Gamma_2,\Gamma_{2\omega},\Gamma_{2\omega^2}$
  of $\twoT$.
\item Leading-order $m_u=m_d$ from T-matrix degeneracy
  $T_{(1,0)}=T_{(0,1)}$; consistent with $m_u/m_d\approx0.46$.
\item Three energy concentration count is exactly $3$: the isometry
  $t\mapsto t+\pi$ combined with $z\mapsto -z$ forces
  $\kappa_{\mathrm{over}}=\kappa_{\mathrm{under}}=0.3991$ at all crossings
  (Remark~\ref{P15:rem:3or6}), so each crossing is a single symmetric
  concentration; higher multiples (6, 9, 12, \ldots) are ruled out.
\item Y-junction / Fermat--Steiner geometry (Proposition~\ref{P15:prop:yjunction}):
  the three crossing midpoints lie at distance exactly $R_0=3$ from the
  origin, forming a regular equilateral triangle.
  The origin is the Fermat--Steiner point with three equal arms of length
  $R_0$; total Y-junction string length $3R_0 = N_3^2 = 9$.
  The missing string-tension term $\sigma\cdot 3R_0$ in the baryon energy
  explains why gradient flows failed to converge (Remark~\ref{P15:rem:yjunction_energy});
  crucially, this term is $C^*$-independent, so the profile saddle
  (determining $C_3^*$ and $\lam_3$) is unchanged.
\item $\vp^9\sqrt{\Ja^{(3)}\Jfour^{(3)}}=3$, equivalent to $\lam_3=\vp^6$
  (Theorems~\ref{P15:thm:norm3}--\ref{P15:thm:equivalence}).
  The condensate volume $V_3=9/(5\vp^{17})$ and saddle value $C_3^*=2.5062$
  are derived consequences (Corollary~\ref{P15:cor:V3}), not independent assumptions.
\item $\lam_3=\vp^6$ from the triple-fusion multiplicity argument
  (Theorem~\ref{P15:thm:sector_assignment}, Corollary~\ref{P15:cor:lambda3_proved}):
  the triple fusion $j_{1/2}^{\otimes 3}=j_{1/2}^{(\times 2)}\oplus j_{3/2}^{(\times 1)}$
  at $k=3$ has a unique multiplicity-1 channel $j=3/2$;
  the topological non-degeneracy of $\pi_3(S^2)=\mathbb{Z}$
  forces $\QH=3$ into this channel;
  the simple-current self-fusion $N_{\mathrm{fus}}=1$ then gives
  $\lam_3=\vp^6$ (three lines).
  All four previously conditional results now follow as proved
  consequences (Corollary~\ref{P15:cor:lambda3_proved}).
\item The $O(\langle\kappa^4\rangle)$ curvature correction is computed
  exactly via Beta-function integration of the BS profile
  (Proposition~\ref{P15:prop:kappa4}): $\delta r_3^{(\kappa^4)}=
  9\langle\kappa^4\rangle/(4 C_3^{*4})=0.000938$, giving
  $C_3^*=2.5062$ (a $0.392\%$ shift from the $O(\langle\kappa^2\rangle)$
  value).
\item The $2I\to 2T$ branching rules are proved by character theory
  (Proposition~\ref{P15:prop:branching}): $j=\tfrac32\to\Gamma_{2\omega}
  \oplus\Gamma_{2\omega^2}$ shows the simple-current sector spans only
  two of the three quark colours, the branching-rule counterpart of
  the sector-assignment proof. The 24-cell geometry of $\twoT$
  (edge-to-circumradius ratio~$1$, vs.\ $1/\vp$ for the 600-cell of
  $\twoI$) explains the integer normalisation
  $\vp^9\sqrt{\Ja^{(3)}\Jfour^{(3)}}=3$ (Remark~\ref{P15:rem:F4}).
\end{enumerate}

%══════════════════════════════════════════════════════════════════════════════
\section{Open Problems}
\label{P15:sec:open}
%══════════════════════════════════════════════════════════════════════════════

\begin{enumerate}[noitemsep,label=(\roman*)]
\item \textbf{[Resolved] Completing the proof of $\lam_3=\vp^6$.}
  Theorem~\ref{P15:thm:sector_assignment} and
  Corollary~\ref{P15:cor:lambda3_proved} close this problem.
  The triple-fusion decomposition
  $j_{1/2}^{\otimes 3}=j_{1/2}^{(\times 2)}\oplus j_{3/2}^{(\times 1)}$
  at $k=3$ (Lemma~\ref{P15:lem:triple_fusion}), combined with
  the topological non-degeneracy $\pi_3(S^2)=\mathbb{Z}$
  (which forbids the multiplicity-2 channel $j=1/2$),
  uniquely assigns $\QH=3$ to the simple-current sector $j=3/2$.
  The self-fusion $N_{\mathrm{fus}}(\tfrac32\otimes\tfrac32)=1$
  then gives $\lam_3=\vp^6$ in three lines
  (Theorem~\ref{P15:thm:lambda3_sc}).
  The normalisation identity, condensate volume, and saddle value
  follow as derived consequences
  (Corollary~\ref{P15:cor:lambda3_proved},
  Theorems~\ref{P15:thm:norm3}--\ref{P15:thm:equivalence},
  Corollary~\ref{P15:cor:V3}).

\item \textbf{[Resolved at leading order] Higher-order curvature
  corrections.}
  Torsion is proved to cancel at leading order, for the
  $\chi$-independent profile $f_0(\rho)$ alone
  (Theorem~\ref{P15:thm:torsion_cancels}).
  The $O(\langle\kappa^4\rangle)$ correction is now computed in closed form
  via Beta-function integration (Proposition~\ref{P15:prop:kappa4}),
  giving $C_3^* = 2.5062$ (shift of $0.392\%$ from the
  $O(\langle\kappa^2\rangle)$ value $2.497$).
  A torsion-coupled correction, structurally analogous to Paper~XIV's
  $f_2(\rho)\cos(2\chi)$ but sourced by $\kappa^2\tau^2$ rather than
  $1/R_0^2$ alone, enters at $O(\rho^2\kappa^2\tau^2)$
  (Remark~\ref{P15:rem:torsion_scope}); numerically
  $\langle\kappa^2\tau^2\rangle\approx 0.0053$, about $4\%$ of
  $\langle\kappa^2\rangle$. This is smaller than but not obviously
  negligible relative to the $O(\langle\kappa^4\rangle)$ term above,
  and its precise coefficient --- together with a structural finding
  that the governing boundary-value problem has complex indicial
  exponents near $\rho=0$, unlike Paper~XIV's torus case --- is
  developed in Paper~XVI~\cite{Paper16}. It does not affect
  $\lam_3=2\vp/r_3$, which is fixed independently of the curvature
  and torsion corrections
  (Theorem~\ref{P15:thm:sector_assignment}).

\item \textbf{Topology-preserving saddle finder.}
  All gradient flows tried on $80^3$ grids --- direct $\Kfb$
  minimisation, fixed-$\lam$ flows at $\lam\in\{\vp^5,\vp^6\}$,
  and projected-gradient flow at fixed $J_4$ --- show $J_4\to 0$
  (topology collapse) without converging to a saddle.
  The collapse is accelerated by the three-lobed energy landscape
  from the crossing-region self-interaction
  (Remark~\ref{P15:rem:no_lambda3}): at the initialisation value
  $C^*=1.5$, the tube radius equals only $2.6$ times the minimum
  inter-tube distance, causing significant inter-tube field overlap.
  Under $\Kfb$ minimisation, $\Kfb/J_4$ forms a $V$-shape (minimum
  $\approx 3.55$, crossing $\vp^6$ twice); under fixed-$\lam$ flows
  it forms an inverted-$U$ (peak $\approx 21\vp^6$ for $\lam=\vp^6$).
  A topology-preserving algorithm --- e.g.\ Battye--Sutcliffe
  simulated annealing~\cite{Battye1998}, constrained flow at fixed
  $J_4$, or an initialisation at $C^*\approx 2.499$
  (separation $4.4$ tube radii) will to confirm
  $C_3^*=2.5062$ and thereby verify the geometric self-consistency
  check~\eqref{P15:eq:V3formula}.
  Since $\lam_3=\vp^6$ is proved analytically, this
  becomes a purely numerical verification rather than a precondition.

\item \textbf{Quark masses from $\twoI\to\twoT$ branching.}
  The generation assignment and T-matrix phases (parts~(a)--(b) below)
  are determined by the branching rules; the absolute mass scale
  (part~(c)) additionally requires the quark Yukawa formula.

  \textbf{(a) Generation-to-irrep assignment.}
  The $\mathrm{SU}(2)$ spin-$j$ representations restrict to $\twoT$ as
  $j=0\to\Gamma_1$, $j=\tfrac12\to\Gamma_2$, $j=1\to\Gamma_3$
  (Proposition~\ref{P15:prop:branching}).
  Each quark generation $g$ carries a combined $(\mathrm{SU}(2)_3$
  level $j_g$, colour $\Gamma_{2\omega^c})$ label, with the
  $\mathbb{Z}_3$ automorphism cycling the colour index independently of
  the generation index. This is the unique assignment consistent with
  three generations, three colours, $\mathbb{Z}_3$ symmetry,
  and $\twoT\subset\twoI$.

  \textbf{(b) T-matrix phases per generation.}
  The $\omega$-twisting of the $\twoT$ fermionic irrep does not change
  the $\mathrm{SU}(2)_3$ modular T-matrix eigenvalue.
  The effective T-phase for generation $g$ is therefore
  $T_g^{\mathrm{quark}} = T_{j_g}^{\mathrm{SU}(2)_3}
  \in\{-\tfrac{3}{40},\,\tfrac{3}{40},\,\tfrac{13}{40}\}$,
  the same values used for leptons in Paper~IV~\cite{Paper4}.
  The inter-generation T-phase ratio
  $m_{g'}/m_g|_T = e^{-2\pi Q_3(T_{g'}-T_g)}$ with $Q_3=3$
  gives $O(1)$ corrections only
  ($m_2/m_1|_T \approx 0.06$ from T-phases alone).
  The large observed hierarchy ($m_c/m_u\approx580$)
  is therefore carried by the $\vp$-tower levels, not the T-phases.

  \textbf{(c) Tower level and absolute mass scale.}
  $C_3^*=2.5062$ (Proposition~\ref{P15:prop:kappa4}) is now known analytically.
  The quark tower level $n_q$ and Yukawa coupling $y_q$ are not
  obtainable by direct substitution into the lepton formula
  ($y_e=e^{-25\pi/6}$ of Paper~IV), because the $\QH=3$ sector
  has a different condensate coupling structure.
  Determining $n_q$ requires identifying how the density-feedback
  functional couples to the trefoil saddle, which is deferred to further work.
\end{enumerate}

\begin{thebibliography}{99}

\bibitem{Paper1}
F.~Manfredi,
\textit{The Density-Feedback Faddeev--Niemi Hopfion:
  Exact Algebraic Predictions for Standard-Model Parameters},
Zenodo (2026).
\href{https://doi.org/10.5281/zenodo.19342027}{doi:10.5281/zenodo.19342027}

\bibitem{Paper2}
F.~Manfredi,
\textit{BPS Structure and Fixed-Point Theorem for the
  Density-Feedback Faddeev--Niemi Hopfion},
Zenodo (2026).
\href{https://doi.org/10.5281/zenodo.19363491}{doi:10.5281/zenodo.19363491}

\bibitem{Paper3}
F.~Manfredi,
\textit{Two Constants from One Knot: The Fine Structure Constant and
  the Linking Scale as Exact Predictions of the Icosahedral WZW Condensate},
Zenodo (2026).
\href{https://doi.org/10.5281/zenodo.19478629}{doi:10.5281/zenodo.19478629}

\bibitem{Paper4}
F.~Manfredi,
\textit{The Hopf Spoke, WZW Fermion Mass Renormalization,
  and Three- and Four-Term Formulas for the Fine Structure Constant},
Zenodo (2026).
\href{https://doi.org/10.5281/zenodo.19504729}{doi:10.5281/zenodo.19504729}

\bibitem{Paper12}
F.~Manfredi,
\textit{The Profile Normalisation Conjecture ---
  Proof via CMB Energy Identity and Two-Route Equivalence},
Zenodo (2026).
\href{https://doi.org/10.5281/zenodo.20210460}{doi:10.5281/zenodo.20210460}

\bibitem{Paper14}
F.~Manfredi,
\textit{The Thick-Torus Profile Correction to the Density-Feedback Hopfion},
Zenodo (2026).
\href{https://doi.org/10.5281/zenodo.20479790}{doi:10.5281/zenodo.20479790}

\bibitem{Paper16}
F.~Manfredi,
\textit{Quark Generation Masses in the Density-Feedback Hopfion:
A Survey of Ruled-Out Mechanisms and the $Q_H=3$ Gradient-Flow Landscape},
Preprint (2026).
\href{https://doi.org/10.5281/zenodo.21012650}{doi:10.5281/zenodo.21012650}

\bibitem{Paper18} F.~Manfredi,
\textit{Preimage Topology and Confinement: Preimage Topology and Confinement: The $Q_H=3$ Sector of the Density-Feedback Hopfion},
Zenodo (2026).
\href{https://doi.org/10.5281/zenodo.21047750}{doi:10.5281/zenodo.21047750}

\bibitem{Faddeev1997}
L.~D. Faddeev and A.~J. Niemi,
\textit{Stable knot-like structures in classical field theory},
Nature \textbf{387}, 58 (1997).
\href{https://doi.org/10.1038/387058a0}{doi:10.1038/387058a0}

\bibitem{Battye1998}
R.~A. Battye and P.~M. Sutcliffe,
\textit{Knots as stable soliton solutions in a three-dimensional classical
field theory},
Phys.\ Rev.\ Lett.\ \textbf{81}, 4798 (1998).
\href{https://doi.org/10.1103/PhysRevLett.81.4798}{doi:10.1103/PhysRevLett.81.4798}

\bibitem{McKay1980}
J.~McKay,
\textit{Graphs, singularities, and finite groups},
Proc.\ Symp.\ Pure Math.\ \textbf{37}, 183--186 (1980).
\href{https://doi.org/10.1090/pspum/037}{doi:10.1090/pspum/037}

%\bibitem{Kasai2026} 
%S.~Kasai et al., 
%\textit{Nonequilibrium dynamics of magnetic hopfions driven by spin-orbit torque, \textit}
%Phys. Rev. B \textbf{113}, 134445 (2026).
%\href{https://doi.org/10.1103/v929-88l7}{doi:10.1103/v929-88l7}

\bibitem{Witten1989}
E.~Witten,
\textit{Quantum field theory and the Jones polynomial},
Commun.\ Math.\ Phys.\ \textbf{121}, 351 (1989).
\href{https://doi.org/10.1007/BF01217730}{doi:10.1007/BF01217730}

\end{thebibliography}
\end{document}